\exercisesection{2}

\begin{enumerate}
\def\labelenumi{\arabic{enumi}.}
\item
  设 K 为交换域，A 为二元不定元上的多项式环 K{[}X, Y{]}，且 \(a_{n}\) 为 A 中的主理想 \((XY^{n})\)；序列 \((a_{n})_{n\geqslant1}\) 与 \(a_{0}=A\) 一起构成 A 上的穷竭且分离滤过。设 b 为 A 的主理想 (X)；证明 A 上的拓扑在 b 上诱导出的拓扑严格粗于由 A 上的滤过在 b 上相应的滤过所定义的拓扑（从而后一个滤过当然不同于由 A 上的滤过诱导的滤过）。
\item
  设 K 为特征 \(\neq2\) 的交换域，A 为二元不定元的形式幂级数环 K{[}{[}X, Y{]}{]}。
\end{enumerate}

\begin{enumerate}
\def\labelenumi{(\alph{enumi})}
\item
  证明在 A 中主理想 \(\mathfrak{p} = (\mathrm{X}^2 -\mathrm{Y}^3)\) 是素理想。（若两个形式幂级数的乘积 \(f(\mathbf{X},\mathbf{Y})g(\mathbf{X},\mathbf{Y})\) 被 \(\mathbf{X}^2 -\mathbf{Y}^3\) 整除，先注意到 \(f(\mathrm{T}^3,\mathrm{T}^2) = 0\) 或 \(g(\mathrm{T}^3,\mathrm{T}^2) = 0\) 在形式幂级数环 \(\mathbf{K}[[\mathrm{T}]]\) 中成立；例如假设 \(f(\mathrm{T}^3,\mathrm{T}^2) = 0\)，先证明 \(f(\mathbf{X},\mathbf{Y}^2)\) 被 \(\mathbf{X} - \mathbf{Y}^3\) 以及 \(\mathbf{X} + \mathbf{Y}^3\) 整除，再考察 \(f(-\mathbf{Y}^3 +\mathbf{X},\mathbf{Y}^2)\)，最后证明 \(f(\mathbf{X},\mathbf{Y}^2)\) 被 \(\mathbf{X}^2 -\mathbf{Y}^6\) 整除。)
\item
  设 \(m\) 为 \(AX + AY\) 这一 \(A\) 的极大理想。证明 \(p\) 关于 \(m\)-进拓扑在 \(A\) 上是闭的（关于此拓扑 \(A\) 是 Hausdorff 且完备的），但在环 \(\text{gr}(A)\) 中，\(\text{gr}(p)\) 不是素理想（\(p\) 赋予由 \(A\) 上滤过诱导的滤过）。
\end{enumerate}

\begin{enumerate}
\def\labelenumi{\arabic{enumi}.}
\setcounter{enumi}{2}
\item
  设 A 为滤过环，E 为有限生成 A-模。证明若 E 赋予由 A 上滤过诱导的滤过，则 gr(E) 是有限生成的 gr(A)-模（参见习题 5 (c)）。~
\item
  给出一个双射 A-线性映射 \(u \colon \mathbf{E} \to \mathbf{F}\) 的例子，其中 \(\mathbf{E}\) 和 \(\mathbf{F}\) 是两个滤过 A-模，\(u\) 与滤过相容，但 \(\operatorname{gr}(u)\) 既非单射也非满射。（取 \(\mathbf{E} = \mathbf{A}\)、\(\mathbf{F} = \mathbf{A}\) 但赋予另一滤过，并取 \(u\) 为恒等映射。）
\item
  设 A 为带滤过 \((\mathfrak{a}_n)_{n\geqslant 0}\) 的交换滤过环，且 \(\mathfrak{a}_0 = \mathbf{A}\)。
\end{enumerate}

\begin{enumerate}
\def\labelenumi{(\alph{enumi})}
\tightlist
\item
  要使环 gr(A) 由次数有界的一族元素生成，充要条件是存在整数 q，使得对所有 \(n\)，\(\mathfrak{a}_n = \mathfrak{a}_{n+1} + \mathfrak{b}_n\)，其中 \(\mathfrak{b}_n\) 是理想 \(\mathfrak{a}_1^{\alpha_1}\mathfrak{a}_2^{\alpha_2}\ldots\mathfrak{a}_q^{\alpha_q}\) 的各项之和：
\end{enumerate}

\[
\alpha_ {i} \geqslant 0 (1 \leqslant i \leqslant q)
\]

\[
\sum_ {i = 1} ^ {q} i \alpha_ {i} = n.
\]

\[
\mathfrak {a} _ {n} = \mathfrak {a} _ {n + k} + \mathfrak {b} _ {n}
\]

\[
k > 0
\]

\begin{enumerate}
\def\labelenumi{(\alph{enumi})}
\setcounter{enumi}{1}
\item
  要使 \(\mathrm{gr}(\mathbf{A})\) 为 Noether 环，充要条件是 \(\mathbf{A} / \mathfrak{a}_1\) 为 Noether 环，(a) 的条件成立，并且对 \(i \leqslant q\)，\(\mathfrak{a}_i / \mathfrak{a}_{i+1}\) 是有限生成的 \((\mathbf{A} / \mathfrak{a}_1)\)-模（使用第 10 节定理 2 的推论 4）。
\item
  设 K 为交换域，A 为二元不定元的多项式环 K{[}X, Y{]}。在单项式 \(X^{m}Y^{n}\) 的集合上定义全序：规定当 \(X^{m}Y^{n} \leqslant X^{p}Y^{q}\)，或当 \(m + n \leqslant p + q\) 且 \(m + n = p + q\) 时，\(m \leqslant p\)。令 \((\mathbf{M}_{k})\) 为按此升序排列的 X, Y 中单项式序列，并令 \(a_{n}\) 为由指标 \(M_{k}\) 的 \(k \geqslant n\) 生成的 A 的理想。证明 \((a_{n})\) 是与 A 的环结构相容的滤过；关于此滤过，gr(A) 有零因子且不是 Noether 环，尽管 A 是 Noether 整环（使用 (b) 中的判据）；特别地，gr( \(a_{1}\) )（在 \(a_{1}\) 上取由 A 上滤过诱导的滤过）不是有限生成的 gr(A)-模，尽管 \(a_{1}\) 是有限生成的 A-模（参见 § 1, 第 2 节，命题 1 的推论）。~
\end{enumerate}

¶ 6. 设 A 为交换环，E 和 F 为两个 A-模，而 (E \(_{n}\) )（分别为 (F \(_{n}\) )）是 E（分别为 F）上由子 A-模组成的穷竭滤过。在张量积 G = E ⊗A F 上，考虑由 G \(_{n}\) = ∑ \(_{i+j=n}\) Im(E \(_{i}\) ⊗A F \(_{j}\) ) 组成的穷竭滤过。

\begin{enumerate}
\def\labelenumi{(\alph{enumi})}
\tightlist
\item
  证明复合典范同态
\end{enumerate}

\[
\begin{array}{r l}&(\mathrm{E} _ {i} / \mathrm{E} _ {i + 1}) \otimes_ {\mathrm{A}} (\mathrm{F} _ {j} / \mathrm{F} _ {j + 1}) \rightarrow\\&\qquad (\mathrm{E} _ {i} \otimes_ {\mathrm{A}} \mathrm{F} _ {j}) / (\operatorname{Im} (\mathrm{E} _ {i} \otimes_ {\mathrm{A}} \mathrm{F} _ {j + 1}) + \operatorname{Im} (\mathrm{E} _ {i + 1} \otimes_ {\mathrm{A}} \mathrm{F} _ {j})) \rightarrow \mathrm{G} _ {i + j} / \mathrm{G} _ {i + j + 1}\end{array}
\]

是一个次数为 0 的分次同态（称为典范同态）\(\operatorname{gr}(\mathbf{E}) \otimes_{\mathbf{A}} \operatorname{gr}(\mathbf{F}) \to \operatorname{gr}(\mathbf{E} \otimes_{\mathbf{A}} \mathbf{F})\) 的限制，该同态为满射。

\begin{enumerate}
\def\labelenumi{(\alph{enumi})}
\setcounter{enumi}{1}
\item
  证明若 \(\operatorname{gr}(\mathbf{E})\) 是平坦 A-模，则当 \(\mathrm{E}_m / \mathrm{E}_n\) 时 A-模 \(m \leqslant n\) 以及当 \(\mathrm{E} / \mathrm{E}_n\) 时 \(n \in \mathbf{Z}\) 都是平坦的。
\item
  以下假设 \(\mathrm{gr}(\mathbf{E})\) 是平坦 A-模，且 \(\mathbf{E}\) 是平坦 A-模（若 \((\mathbf{E}_n)\) 是离散滤过，第二个假设由第一个假设推出）。于是 \(\mathrm{H}_{ij} = \mathrm{E}_i \otimes_{\mathbf{A}} \mathrm{F}_j\) 被典范地识别为 \(\mathbf{G} = \mathbf{E} \otimes_{\mathbf{A}} \mathbf{F}\) 的子模（第一章，第 2 节，第 5 节，命题 4）。证明对于 \(\mathbf{R}\) 的任意两个有限子集 \(\mathbf{S}\)、\(\mathbf{Z} \times \mathbf{Z}\)，
\end{enumerate}

\[
\left(\sum_ {(i, j) \in \mathbf {R}} \mathrm{H} _ {i j}\right) \cap \left(\sum_ {(h, k) \in \mathbf {S}} \mathrm{H} _ {h k}\right) = \mathrm{H} _ {\sup (i, h), \sup (j, k)}\tag{*}
\]

其中第二个和式中的 \((i,j,h,k)\) 遍历 \(R \times S\)。（若 R 和 S 各自只含一个元素，使用第一章，第 2 节，第 6 节，命题 7。若 p 是 R 或 S 的元素中出现的指标 i 和 h 的最大值，q 是其中出现的指标 j 和 k 的最小值，考察 (*) 两边在典范同态

\[
\mathbf {E} \otimes_ {\mathbf {A}} \mathbf {F} _ {q} \rightarrow (\mathrm{E} / \mathrm{E} _ {p}) \otimes_ {\mathbf {A}} \mathbf {F} _ {q}
\]

下的像，并对 \(\operatorname{Card}(\mathbf{R}) + \operatorname{Card}(\mathbf{S})\) 作归纳。）

\begin{enumerate}
\def\labelenumi{(\alph{enumi})}
\setcounter{enumi}{3}
\tightlist
\item
  由 (c) 推出 (a) 中定义的典范同态于是为双射。
\end{enumerate}

\begin{enumerate}
\def\labelenumi{\arabic{enumi}.}
\setcounter{enumi}{6}
\tightlist
\item
  设 A 为交换环，E 为 A-模，F 为 E 的子模，B 为外代数 \(\wedge\) (E)，而 \(\Im\) 为 B 中由 F 的元素在 B 中的典范像生成的双边理想。令 \(\operatorname{gr}^{\Im}(B)\) 为 B 按 \(\Im\)-进滤过所得的伴随分次环。
\end{enumerate}

\begin{enumerate}
\def\labelenumi{(\alph{enumi})}
\item
   定义一个满射的典范（非分次）A-代数同态 \((\wedge (\mathbf{F})) \otimes_{\mathbf{A}}^{g}(\wedge (\mathbf{E} / \mathbf{F})) \to \mathrm{gr}^{\Im} (\mathbf{B})\)（这里指分次代数的斜张量积（代数，第 III 章）。
\item
  证明若 \(\mathbf{F}\) 在 \(\mathbf{E}\) 中有补子模，则 (a) 中定义的同态是同构。
\item
   取 \(\mathbf{A} = \mathbf{Z}\)、\(\mathbf{E} = \mathbf{Z} / 4\mathbf{Z}\)、\(\mathbf{F} = 2\mathbf{Z} / 4\mathbf{Z}\)；证明此时 (a) 中定义的同态不是单射。
\end{enumerate}

\begin{enumerate}
\def\labelenumi{\arabic{enumi}.}
\setcounter{enumi}{7}
\item
  设 A 为滤过环，\((\mathbf{A}_{n})\) 为其滤过，E 为滤过 A-模，\((\mathbf{E}_{n})\) 为其滤过。证明若 \(A_{0} = A\) 且 \(E_{0} = E\)，则从 \((a, x) \mapsto ax\) 到 E 的映射 \(A \times E\) 关于由这些滤过定义的拓扑是一致连续的。
\item
  给出两个滤过环 A、B 的例子，它们的滤过是穷竭且分离的，并给出一个与滤过相容的非满射同态 \(u: A \to B\)，使得 \(\text{gr}(u)\) 为双射。由此在不假设环 A 完备时，给出命题 12 的推论 1 的反例；并在不再假设 E 有限生成时，给出命题 13 的反例（使用代数，第 VIII 章，第 7 节，习题 3(b)）。
\item
  设 A 为 Noether 环，\(\sigma\) 为 A 的自同构。证明代数，第 IV 章，第 5 节，习题 10(b) 中定义的环 E 是（左或右）Noether 环。
\item
  证明若 E 是交换域上维数 \(\geqslant2\) 的向量空间，则 E 的张量代数是既非左 Noether 环也非右 Noether 环的环（若 a、b 是 E 中两个线性无关的向量，考察由 \(a^{n}b^{n}\) 时的元素 \(n\geqslant1\) 生成的左理想（或右理想））。
\end{enumerate}

 ¶ 12. 设 K 为交换域，A 为任意无限族不定元上的多项式环 \(K[X_{i}]_{i\in I}\)，m 为 A 中由 \(X_{i}\) 生成的（极大）理想。若令 \(A_{i}=A/m^{i+1}\)，则 \(A_{i}\) 以及当 \(h_{ij}:A/m^{j+1}\to A/m^{i+1}\) 时的典范同态 \(i\leqslant j\) 满足命题

14；环 \(\lim A_{i}\) 是完备化 \(\hat{A}\)（关于 m-进拓扑的 \(A\)），且典范同态 \(\hat{A} \to A_{i}\) 的核等于 \((m^{i})^{\wedge}\)，即 \(m^{i}\) 在 \(\hat{A}\) 中的闭包。

\begin{enumerate}
\def\labelenumi{(\alph{enumi})}
\item
  证明 \(\hat{\mathbf{A}}\) 被典范地识别为关于 \(\mathbf{X}_{\iota}\) 的形式幂级数环，其中每个幂级数在给定次数下只有有限项（代数，第 IV 章，第 5 节，习题 1）。
\item
  从现在起取 \(\mathbf{I} = \mathbf{N}\)。证明 \(\hat{\mathfrak{m}} \neq \hat{\mathbf{A}}. \mathfrak{m}\)（考察形式幂级数 \(\sum_{n=1}^{\infty} \mathbf{X}_n^n\)）。
\item
  设 K 为有限域。证明 \((\hat{\mathfrak{m}})^{2} \neq (\mathfrak{m}^{2})^{\bullet}\)。（先证明如下结果：对每个整数 k \textgreater{} 0，存在整数 \(n_{k}\)，并且对每个整数 \(n \geqslant n_{k}\)，存在一个以 K 为系数、在 \(F_{n}\) 个不定元中次数为 n 的齐次多项式 \(n^{2}\)，使得 \(F_{n}\) 不可能是形如 \(P_{1}Q_{1} + \cdots + P_{k}Q_{k}\) 的任一多项式的 n 次项之和，其中 \(P_{i}\) 和 \(Q_{i}\) 是同一组 \(n^{2}\) 个不定元中的无常数项多项式。）
\end{enumerate}

推出 \(\hat{A}\) 关于 \(\hat{m}\)-进拓扑不是完备的。

\begin{enumerate}
\def\labelenumi{\arabic{enumi}.}
\setcounter{enumi}{12}
\tightlist
\item
  设 K 为域，\(A = K[[X]]\) 为形式幂级数环，m 为其极大理想，故 A 关于 m-进拓扑是 Hausdorff 且完备的（第 6 节，命题 6 的推论）。在加法群 A 上考虑滤过 \((\mathrm{E}_{n})\)，其中 \(E_{0} = A\)，而 \(E_{n}\) 是 \(m^{n}\) 与环 K{[}X{]} 的交；此滤过是穷竭且分离的，它在 A 上定义的拓扑 T 与加法群结构相容并且比 m-进拓扑更细，但 A 关于拓扑 T 不是完备群（考察多项式序列 \((1 - X^{n})/(1 - X)\)）。
\end{enumerate}

14 设 A 为环（不必交换），E 为左 A-模。若 E 上的拓扑在平移下不变且 0 有一个由 E 的子模组成的邻域基本系，则称该拓扑为线性拓扑；此时称 E 带线性拓扑。线性拓扑与 E 的加法群结构相容，并且当 A 取离散拓扑时在 E 上定义拓扑 A-模结构。任意 A-模上的离散拓扑和最粗拓扑都是线性拓扑。

\begin{enumerate}
\def\labelenumi{(\alph{enumi})}
\item
  若 E 是带线性拓扑的模，F 是 E 的子模，则 F 上的诱导拓扑以及 E 关于 F 的商拓扑都是线性拓扑。若 \((\mathrm{E}_{\alpha}, f_{\alpha\beta})\) 是带线性拓扑的 A-模的逆系统，其中 \(f_{\alpha\beta}\) 是连续线性映射，则拓扑 A-模 \(E = \lim_{\leftarrow} E_{\alpha}\) 带线性拓扑。
\item
  设 E 为带线性拓扑的 A-模。E 中存在一个由开（也是闭的）子模组成的 0 的邻域基本系 \((\mathrm{V}_{\lambda})_{\lambda \in L}\)；若 \(\phi_{\lambda\mu}: E/V_{\mu} \to E/V_{\lambda}\) 为典范映射且 \(V_{\lambda} \supset V_{\mu}\)，则族 \((\mathrm{E}/\mathrm{V}_{\lambda}, \phi_{\lambda\mu})\) 是离散 A-模的逆系统；拓扑 A-模 \(\tilde{E} = \lim_{\leftarrow} E/V_{\lambda}\) 被识别为 E 的 Hausdorff 完备化，因而其带线性拓扑（一般拓扑，第 III 章，第 7 节，第 3 节）。
\item
  Hausdorff 的带线性拓扑 A-模 E 若为 Artin 模，或 E 的子模 \(\neq0\) 的集合中存在最小元，则称其为离散的。
\end{enumerate}

\begin{enumerate}
\def\labelenumi{\arabic{enumi}.}
\setcounter{enumi}{14}
\tightlist
\item
  \begin{enumerate}
  \def\labelenumii{(\alph{enumii})}
  \tightlist
  \item
    设 E 为带线性拓扑的 A-模（习题 14）。由 E 中线性簇组成的滤基具有聚点的充要条件是存在一个包含它的、由仿射线性簇组成的收敛滤基 \(B' \supset B\)。若 E 是 Hausdorff 的，且 E 上每个由仿射线性簇组成的滤基都至少有一个聚点，则称 E 是线性紧的。每个 Artin 模在离散拓扑下都是线性紧的。Hausdorff 的带线性拓扑模 F 的每个线性紧子模在 F 中都是闭的。
  \end{enumerate}
\end{enumerate}

\begin{enumerate}
\def\labelenumi{(\alph{enumi})}
\setcounter{enumi}{1}
\item
  若 E 是线性紧 A-模，且 u 是从 E 到 Hausdorff 的带线性拓扑 A-模 F 的连续线性映射，则 \(u(\mathbf{F})\) 是 F 的线性紧子模。
\item
  设 E 为 Hausdorff 的带线性拓扑 A-模，F 为 E 的闭子模。E 线性紧的充要条件是 F 和 E/F 都线性紧。
\item
  线性紧模的任意乘积都是线性紧的（考察在线性仿射簇滤基中极大的滤基。
\item
  设 \((\mathrm{E}_{\alpha}, f_{\alpha \beta})\) 是关于有向指标集的带线性拓扑模的逆系统；假设 \(f_{\alpha \beta}\) 是连续线性映射，且当 \(\alpha \leqslant \beta\) 时，\(f_{\alpha \beta}^{-1}(0)\) 是 \(\mathbf{E}_{\beta}\) 的线性紧子模。令 \(\mathbf{E} = \varprojlim \mathbf{E}_{\alpha}\)，令 \(f_{\alpha}\) 为典范映射 \(\mathbf{E} \to \mathbf{E}_{\alpha}\)；证明对所有 \(\alpha, f_{\alpha}(\mathbf{E}) = \bigcap_{\alpha \leqslant \beta} f_{\alpha \beta}(\mathbf{E}_{\beta})\)（特别地，若 \(f_{\alpha \beta}\) 为满射，则 \(f_{\alpha}\) 也是满射），且 \(f_{\alpha}^{-1}(0)\) 是线性紧的（使用 (d) 以及一般拓扑，第 I 章，附录，第 2 节，定理 1）。
\end{enumerate}

¶ 16. (a) 设 E 为 Hausdorff 的带线性拓扑 A-模。证明下列性质等价：

(α) E 是线性紧的（习题 15）。

(β) 对每个连续线性映射 \(u\)，它从 \(E\) 映到 Hausdorff 的带线性拓扑 A-模 \(F\)，\(u(E)\) 都是 \(F\) 的闭子模。

\begin{enumerate}
\def\labelenumi{(\Alph{enumi})}
\setcounter{enumi}{24}
\tightlist
\item[(γ)]
  对 E 上每个比给定拓扑粗的线性拓扑（无论是否 Hausdorff），E 都是完备的。
\end{enumerate}

(δ) E 是完备的，且 E 中存在一个由子模组成的 0 的开邻域基本系 \((\mathbf{U}_{\lambda})\)，使得 \(E/U_{\lambda}\) 是离散的线性紧 A-模（参见 § 3，习题 5）。~

（为看出 \((\gamma)\) 蕴含 \((\delta)\)，取 \(F\) 为 \(E\) 的一个开子模，以及由仿射线性簇组成的滤基 \(\mathfrak{B}\) 在 \(E / F\) 上。对所有 \(V \in \mathfrak{B}\)，令 \(M_V\) 为 V 在 E/F 中的方向子模在 E 中的逆像；在 E 上考虑以 \(M_{V}\) 为 0 的邻域基本系的线性拓扑。）

\begin{enumerate}
\def\labelenumi{(\alph{enumi})}
\setcounter{enumi}{1}
\item
  设 E 为 Hausdorff 的带线性拓扑 A-模，M 为 E 的线性紧子模。证明对于 E 的每个闭子模 F，\(M + F\) 在 E 中是闭的（考察 M 在 E/F 中的像）。
\item
  设 E 为线性紧 A-模，F 为 Hausdorff 的带线性拓扑 A-模。证明 \(E \times F\) 的每个闭子模 M 在 F 上的投影都是闭的。得到逆命题（参见一般拓扑，第 I 章，第 10 节，第 2 节）。~
\item
  设 E 为线性紧 A-模，u 为从 E 到 Hausdorff 的带线性拓扑 A-模 F 的连续线性映射。证明对于 E 上每个由仿射线性簇组成的滤基 B，B 的聚点集在 u 下的像就是 \(u(\mathfrak{B})\) 的聚点集。特别地，对于 E 的每个闭子模 M，都有 \(\bigcap_{N\in\mathfrak{B}}(M+N)=M+\bigcap_{N\in\mathfrak{B}}N\)。
\end{enumerate}

\begin{enumerate}
\def\labelenumi{\arabic{enumi}.}
\setcounter{enumi}{16}
\tightlist
\item
  若 A-模 E 上的线性拓扑 \(\mathcal{T}\) 是 Hausdorff 的，且不存在严格粗于 \(\mathcal{T}\) 的 Hausdorff 线性拓扑，则称其为极小的。
\end{enumerate}

\begin{enumerate}
\def\labelenumi{(\alph{enumi})}
\item
  Hausdorff 线性拓扑 T 在 E 上为极小的充要条件是：E 上每个由仿射线性簇组成且只有一个聚点的滤基 B 都收敛到该点。（为看出条件是必要的，注意当 M 遍历 B、V 遍历由子模组成的 0 的邻域基本系时，M + V 构成一个与 B 有相同聚点的滤基。为看出条件是充分的，注意由开子模组成且交约化为 0 的滤基，是比 T 粗的某个 Hausdorff 线性拓扑中 0 的邻域基本系。）
\item
  A-模 E 上的离散拓扑为极小的充要条件是 E 为 Artin 模。
\item
  若 \(\mathcal{T}\) 为极小拓扑，则 \(\mathcal{T}\) 在 \(\mathbf{E}\) 的任意闭子模上诱导的拓扑也是极小的。
\end{enumerate}

¶ 18. 在 A-模 E 中，若子模 M ≠ E 且 E/M 的非零子模集合中存在最小元，则称 M 为受保护的。

\begin{enumerate}
\def\labelenumi{(\alph{enumi})}
\item
  证明 E 的每个子模 \(N \neq E\) 都是若干受保护子模的交（对于所有 \(x \notin N\)，考察 E 中包含 N 但不包含 x 的子模集合中的极大元）。
\item
  设 \(\mathcal{T}\) 是 E 上的 Hausdorff 线性拓扑，令 \(\mathcal{T}^*\) 为如下线性拓扑：其 0 的邻域基本系是由 E 中关于 \(\mathcal{T}\) 开且受保护的子模生成的滤基。证明 \(\mathcal{T}^*\) 是 Hausdorff 的，且关于 \(\mathcal{T}\) 闭的每个子模关于 \(\mathcal{T}^*\) 也闭（注意关于 \(\mathcal{T}\) 闭的每个子模都是关于 \(\mathcal{T}\) 开的子模的交）。推出若 E 关于 \(\mathcal{T}^*\) 完备，则 E 关于 \(\mathcal{T}\) 完备（一般拓扑，第 III 章，第 3 节，第 5 节，命题 9）。
\item
  假设 \(\mathcal{T}\) 是线性紧的。证明 \(\mathcal{T}^*\) 是线性紧的，并且是比 \(\mathcal{T}\) 粗的 Hausdorff 线性拓扑中最粗的；特别地，\(\mathcal{T}^*\) 是极小线性拓扑（习题 17）。（令 \(\mathfrak{B}\) 为由关于 \(\mathcal{T}\) 闭的子模组成且交为 0 的滤基；使用习题 16 (d)，证明对于关于 \(\mathcal{T}\) 开且受保护的每个子模 U，存在 \(\mathbf{M} \in \mathfrak{B}\) 使得 \(\mathbf{M} \subset \mathbf{U}\)。）
\item
  设 \(\mathbf{F}\) 为 Hausdorff 的带线性拓扑 A-模，\(\mathcal{T}_1\) 为其拓扑；证明若 \(u: \mathrm{E} \to \mathrm{F}\) 关于拓扑 \(\mathcal{T}\) 和 \(\mathcal{T}_1\) 是连续线性映射，则 \(u\) 关于拓扑 \(\mathcal{T}^*\) 和 \(\mathcal{T}_1^*\) 也连续。
\end{enumerate}

¶ 19. (a) 设 E 为线性紧 A-模。证明下列条件等价：

(α) 对从 E 到 Hausdorff 的带线性拓扑 A-模 F 的每个连续线性映射 u，u 都是从 E 到 F 的严格态射（一般拓扑，第 III 章，第 2 节，第 8 节）。

(β) 对 \(\mathbf{F}\) 的每个闭子模 \(\mathbf{E}\)，\(\mathbf{E} / \mathbf{F}\) 上的商拓扑都是极小拓扑（习题 17）。

(γ) E 是完备的，且 E 中存在一个由子模组成的 0 的开邻域基本系 \((\mathbf{U}_{\lambda})\)，使得 \(E/U_{\lambda}\) 是 Artin 模。

（为看出 \((\gamma)\) 蕴含 \((\beta)\)，将其化约到 \(F = 0\) 的情形，并使用习题 17 和 18。）

若 E 满足等价条件 \((\alpha)\)、\((\beta)\) 和 \((\gamma)\)，则称 E 为严格线性紧的。

\begin{enumerate}
\def\labelenumi{(\alph{enumi})}
\setcounter{enumi}{1}
\item
  设 E 为 Hausdorff 的带线性拓扑 A-模，F 为 E 的闭子模。E 严格线性紧的充要条件是 F 和 E/F 都严格线性紧。
\item
  严格线性紧模的每个逆极限都是严格线性紧的。
\item
  设 E 为 Hausdorff 的带线性拓扑 A-模，u 为从 E 到严格线性紧 A-模 F 的线性映射。证明若 u 的图像在 \(E \times F\) 中闭，则 u 连续（使用习题 16 (c) 以及如下事实：若 M 在 E 中闭且 E/M 是 Artin 模，则 M 在 E 中开）。
\end{enumerate}

¶ 20. (a) 证明离散的线性紧模不可能是无穷多个不约化为 0 的子模的直和。

\begin{enumerate}
\def\labelenumi{(\alph{enumi})}
\setcounter{enumi}{1}
\item
  给出一个不是线性紧的极小线性拓扑（习题 17）的例子（考察无穷个两两不同构的简单模的直和）。
\item
  设 E 为线性紧模，且存在一个简单子模族，其和在 E 中稠密。证明：(1) E 严格线性紧（将 (a) 应用于 E 对一个开子模的商）；(2) E 同构于一族离散简单模的（拓扑）乘积。（考察 E 的极大开子模集合 O，使得对于每个有限
\end{enumerate}

\[
\mathrm{E} / \left(\bigcap_ {k = 1} ^ {n} \mathrm{G} _ {k}\right)
\]

 个由序列 \((\mathbf{G}_k)_{1\leqslant k\leqslant n}\) 中的 \(n\) 个属于 \(\mathfrak{O}\) 的互异元素组成的序列，\(\operatorname{E} / \left(\bigcap_{k=1}^{\infty}\mathbf{G}_k\right)\) 是 \(n\) 个简单子模的直和。证明存在一个（关于包含关系）极大集合 \(\mathfrak{O}\)，且属于此集合的所有子模 \(\mathbf{G}\) 的交约化为 0；使用如下事实：\(\{0\}\) 是极大开子模的交，并且对于每个极大开子模 \(\mathbf{G}\)，它是 \(\mathbf{E}\) 的子模，\(\mathbf{E}\) 至少有一个不包含于 \(\mathbf{G}\) 的简单子模。最后使用 \(\mathbf{E}\) 严格线性紧这一事实得出结论。）

\begin{enumerate}
\def\labelenumi{(\alph{enumi})}
\setcounter{enumi}{3}
\tightlist
\item
  由 (c) 推出每个域 \(\mathbf{K}\) 上的线性紧向量空间都是严格线性紧的，并同构于乘积 \(\mathbf{K}_s^{\mathrm{I}}\)。
\end{enumerate}

¶ 21. 回忆一下，在环 A 上，若一个拓扑是 A-模 \(\mathbf{A}_s\) 上的线性拓扑且与 A 的环结构相容，则称其为（左）线性的（第二章，第 2 节，习题 16）。若拓扑环 A 的 \(\mathbf{A}_s\) 是线性紧（分别严格线性紧）A-模，则称 A 为（左）线性紧（分别（左）严格线性紧）的。

\begin{enumerate}
\def\labelenumi{(\alph{enumi})}
\item
  证明若 A 关于拓扑 T 是左线性紧的，则习题 18 (b) 中定义的拓扑 \(T^{*}\) 也与 A 的环结构相容（使用习题 18 (d)）。
\item
  假设 \(\mathbf{A}\) 是交换的；令 \(u \neq 0\) 为 \(\mathbf{A} / \Re\) 的幂等元，其中 \(\Re\) 是 \(\mathbf{A}\) 的 Jacobson 根。证明若 \(\mathbf{A}\) 取离散拓扑且为线性紧的，则存在唯一的幂等元 \(e \in \mathbf{A}\)，其在 \(\mathbf{A} / \Re\) 中的像等于 \(u\)。（在线性仿射簇 \(x + \mathfrak{b}\) 中，其中 \(\mathfrak{b} \subset \Re\) 且 \(x^2 - x \in \mathfrak{b}\)，且该簇包含于类 \(u\)，证明存在最小元 \(e + \mathfrak{a}\)。使用代数，第 VIII 章，第 6 节，习题 10(a)，推出 \(e^2 = e\) 且 \(\mathfrak{a} = 0\)。考察 \(e^2 - e = r\)，证明 \(r \in \mathrm{Ar}^2\)，并证明 \(e\) 的唯一性：若 \(e_1, e_2\) 是 \(\mathbf{A}\) 的两个幂等元且 \(e_1 e_2 \in \Re\)，则 \(e_1 e_2 = 0\)。）
\item
  证明每个取离散拓扑且线性紧的交换环 A 都是有限个局部环的直积（这些局部环也取离散拓扑且线性紧）。（先考虑 \(\Re = 0\) 的情形，使用习题 15 (b) 以及第二章，第 1 节，第 2 节，命题 5；然后将 (b) 应用于 A/ \(\Re\) 的幂等元。）
\item
  证明每个线性紧（分别严格线性紧）交换环 A 都是线性紧（分别严格线性紧）局部环族的乘积。（令 \((\mathfrak{m}_{\lambda})\) 为 A 的开极大理想族；对于每个开理想 \(a\)，它不包含 \(\mathfrak{m}_{\lambda}\)，令 \(\mathrm{A}_{\lambda}^{\prime} \times \mathrm{A}_{\lambda}^{\prime \prime}\) 为 A/a 作为两个环的直积分解，使 \(\mathrm{A}_{\lambda}^{\prime}\) 是 (c) 中定义的、极大理想为 \((\mathfrak{m}_{\lambda} + \mathfrak{a}) / a\) 的局部环；令 \(e_{\lambda}^{\prime}(a)\) 为 \(\mathrm{A}_{\lambda}^{\prime}\) 的单位元，并将其视为 A 中模 \(a\) 的类；\(e_{\lambda}^{\prime}(a)\) 构成一个在 A 中收敛到幂等元 \(e_{\lambda}\) 的滤基，且 \(Ae_{\lambda}\) 是 A 的闭理想，\(e_{\lambda}\) 是其单位元。证明环 \(Ae_{\lambda}\) 是局部环，且 A 同构于各个 \(Ae_{\lambda}\) 的乘积。）
\end{enumerate}

¶ 22. (a) 设 A 为局部环，m 为其极大理想。证明若关于 A 上的线性拓扑 T，A 是严格线性紧的，则 T 比 m-进拓扑粗。

\begin{enumerate}
\def\labelenumi{(\alph{enumi})}
\setcounter{enumi}{1}
\item
  设 A 为交换环，m 为 A 的理想。A 关于 m-进拓扑严格线性紧的充要条件是：A 关于此拓扑 Hausdorff 且完备，A/m 是 Artin 环，且 \(m/m^{2}\) 是一个有限长度的 (A/m)-模 *（换言之，A 是完备的 Noether 半局部环）*。
\item
  设 I 为无限指标集，K 为有限域，\(A = K[[X_{t}]]_{t \in I}\) 为关于不定元族 \((\mathbf{X}_{t})_{t \in I}\) 的形式幂级数代数（代数，第 IV 章，第 5 节，习题 1）；A 是局部环。作为 K 上的向量空间，A 可识别为乘积空间 \(\mathbf{K}^{\mathbf{N}^{(1)}}\)；若 K 取离散拓扑而 A 取乘积拓扑 T，证明 T 是环 A 上的线性拓扑，且 A 关于此拓扑严格线性紧；若 m 是 A 的极大理想，仿照习题 12 (c) 证明 A 关于 m-进拓扑不是完备的。
\end{enumerate}

¶ 23. 设 A 为带线性拓扑 T 的交换环。

\begin{enumerate}
\def\labelenumi{(\alph{enumi})}
\item
  证明 \(\mathfrak{a}\) 的如下理想 \(\mathbf{A}\)：它们关于 \(\mathcal{T}\) 是闭的且 \(\mathbf{A} / \mathfrak{a}\) 是 Artin 环，构成 \(\mathcal{T}^{(c)}\) 上某个线性拓扑 \(\mathbf{A}\) 的 0 的邻域基本系，而该拓扑比 \(\mathcal{T}\) 粗（使用代数，第 VIII 章，第 2 节，习题 12）；然后 \((\mathcal{T}^{(c)})^{(c)} = \mathcal{T}^{(c)}\)。令 \(\mathbf{B}\) 为 \(\mathbf{A}\) 关于拓扑 \(\mathcal{T}^{(c)}\) 的 Hausdorff 完备环；证明 \(\mathbf{B}\) 是严格线性紧的；称 \(\mathbf{B}\) 为与 \(\mathbf{A}\) 相伴的严格线性紧环。
\item
  令 \(\mathcal{T}_{\omega}(\mathbf{A})\) 为 \(\mathbf{A}\) 上的离散拓扑，\(\mathcal{T}_c(\mathbf{A})\) 为拓扑 \((\mathcal{T}_{\omega}(\mathbf{A}))^{(c)}\)。拓扑 \(\mathcal{T}^{(c)}\) 在 \(\mathbf{A}\) 上为 Hausdorff 的充要条件是 \(\mathcal{T}\) 为 Hausdorff，且对于每个开理想 \(\mathfrak{b}\)，它是 \(\mathbf{A}\) 中关于 \(\mathcal{T}\) 开的理想，拓扑 \(\mathcal{T}_c(\mathbf{A}/\mathfrak{b})\) 为 Hausdorff。* 若 \(\mathbf{A}\) 是高度 1 的非离散赋值的环（第 VI 章），则 \(\mathcal{T}_c(\mathbf{A})\) 不是 Hausdorff。*
\item
  假设 A 关于 \(\mathcal{T}\) 是线性紧的；则 A 关于 \(\mathcal{T}^{(c)}\) 是完备的（但一般不是 Hausdorff 的）。\(\mathcal{T}^{(c)}\) 为 Hausdorff 的充要条件是 A 上存在一个严格线性紧且比 \(\mathcal{T}\) 粗的线性拓扑；此时 \(\mathcal{T}^{(c)}\) 是具有这些性质的唯一拓扑。
\item
  设 E 为带线性拓扑且严格线性紧的 A-模。证明若 A 取拓扑 \(\mathcal{T}_{c}(A)\)，则 E 是拓扑环（注意若 F 是 Artin A-模，则对所有 \(x \in F\)，x 的湮灭子 a 使得 A/a 是 Artin 环）。推出若 B 是 A 关于拓扑 \(\mathcal{T}_{c}(A)\) 的 Hausdorff 完备化，则 E 可看作拓扑 B-模；环 B 同构于严格线性紧局部环的乘积
\end{enumerate}

\(B_{\lambda}\)（习题 21 (d)），证明 E 同构于严格线性紧子模 \(E_{\lambda}\) 的乘积，其中 \(E_{\lambda}\) 被指标 \(B_{\mu}\) 的 \(\mu \neq \lambda\) 湮灭，因而可看作拓扑 \(B_{\lambda}\)-模（若 \(e_{\lambda}\) 是 B 的单位元，取 \(E = e_{\lambda}.E\)）。

\begin{enumerate}
\def\labelenumi{\arabic{enumi}.}
\setcounter{enumi}{23}
\tightlist
\item
  在交换环 A 上，令 \(\mathcal{T}_{m}(A)\)~ 表示如下线性拓扑：其 0 的邻域基本系由有限个极大理想的幂的乘积（等于交）组成（参见第 13 节，命题 17）；若 A 中任意有限个理想 \(\neq0\) 的交都是 \(\neq0\)，令 \(\mathcal{T}_{u}(A)\) 表示如下线性拓扑：其 0 的邻域基本系由 A 的所有理想 \(\neq\{0\}\) 组成。
\end{enumerate}

\begin{enumerate}
\def\labelenumi{(\alph{enumi})}
\item
  证明 \(\mathcal{T}_c(\mathbf{A})\)（习题 23）比 \(\mathcal{T}_m(\mathbf{A})\) 粗（注意 Artin 环的 Jacobson 根是幂零的）；给出 \(\mathcal{T}_c(\mathbf{A}) \neq \mathcal{T}_m(\mathbf{A})\) 的例子（参见习题 22 (c)）。~* 若 A 是高度 1 的非离散赋值的环（第 VI 章），则 \(\mathcal{T}_c(\mathbf{A}) = \mathcal{T}_m(\mathbf{A})\)，而 \(\mathcal{T}_u(\mathbf{A}) \neq \mathcal{T}_m(\mathbf{A})\)。*
\item
  拓扑 \(\mathcal{T}_m(A)\) 是 \(A\) 上满足以下条件的最粗线性拓扑：\(A\) 的极大理想都是闭的，且每个开理想的每个幂都是开理想（注意在 \(A\) 上使极大理想 \(m\) 闭的线性拓扑下，\(m\) 必然是开的）。
\item
  若 A 是 Noether 环，则 \(\mathcal{T}_{m}(\mathbf{A}) = \mathcal{T}_{c}(\mathbf{A})\)（注意若 \(m\) 是 A 的极大理想，则对每个整数 \(\mathbf{A} / m^{k}\)，\(k \geqslant 1\) 是 Artin 环，并注意 \(\mathfrak{m}^h / \mathfrak{m}^{h+1}\) 是必然具有有限长度的 (A/m)-模）。给出一个满足 \(\mathcal{T}_{m}(\mathbf{A}) \neq \mathcal{T}_{u}(\mathbf{A})\) 的 Noether 局部环 A 的例子（考察域上多项式环的一个局部环）。
\end{enumerate}

\begin{enumerate}
\def\labelenumi{\arabic{enumi}.}
\setcounter{enumi}{24}
\item
  设 A 为拓扑环，E 为有限生成左 A-模。证明 E 上存在一个拓扑（称为典范拓扑），该拓扑与其 A-模结构相容且比其他所有拓扑都细（将 E 写成 \(A_{s}^{n}/R\) 的形式，在 \(A_{s}^{n}\) 上取乘积拓扑，在 E 上取商拓扑）。若 \(E'\) 是拓扑 A-模且 \(u: E \to E'\) 是 A-线性映射，证明 u 关于 E 上的典范拓扑是连续的。若进一步 \(E'\) 有限生成且带典范拓扑，则 u 是严格态射。若 A 上的拓扑由滤过（\(A_{n}\)）定义，证明 E 上的典范拓扑由滤过（\(A_{n}E\)）定义。
\item
  设 A 为交换环，m 为 A 的理想，\((a_{\lambda})_{\lambda \in L}\) 为 m 的一个生成系，\(\hat{A}\) 为 A 关于 m-进拓扑的 Hausdorff 完备化。
\end{enumerate}

 证明若 \(\mathbf{A}'\) 表示关于不定元族 \((\mathrm{T}_{\lambda})_{\lambda \in \mathbf{L}}\) 的形式幂级数环，且每个给定次数只有有限项（习题 12），则 \(\hat{\mathbf{A}}\) 同构于 \(\mathbf{A}'\) 对理想的闭包 \(\mathfrak{b}\) 的商，其中该理想是 \(\mathbf{A}'\) 中由 \(T_{\lambda} - a_{\lambda}\) 生成的理想（使用第 8 节定理 1）。特别地，p-进整数环 \(Z_{p}\) 同构于商环 \(\mathbf{Z}[[T]]/(T-p)\)。

¶ 27. (a) 设 A 为带线性拓扑的交换环。对于 A 的每个乘性子集 S，令 \(\mathrm{A}\{\mathrm{S}^{-1}\}\) 表示环 \(\mathrm{S}^{-1}\mathrm{A}\) 的 Hausdorff 完备化，其拓扑的 0 的邻域基本系由理想 \(\mathrm{S}^{-1}\mathrm{U}_{\lambda}\) 组成，其中 \((\mathrm{U}_{\lambda})\) 是 A 中由 A 的理想组成的 0 的邻域基本系。证明 \(\mathrm{A}\{\mathrm{S}^{-1}\}\) 典范同构于环 \(\mathrm{S}_{\lambda}^{-1}(\mathrm{A}/\mathrm{U}_{\lambda})\) 的逆极限，其中 \(\mathrm{S}_{\lambda}\) 是 S 在 \(\mathrm{A}/\mathrm{U}_{\lambda}\) 中的典范像。若 \(\hat{\mathrm{A}}\) 是 A 的 Hausdorff 完备化，则 \(\mathrm{A}\{\mathrm{S}^{-1}\}\) 典范同构于 \(\hat{\mathrm{A}}\{\mathrm{S}'^{-1}\}\)，其中 \(\mathrm{S}'\) 是 S 在 \(\hat{\mathrm{A}}\) 中的典范像。

\begin{enumerate}
\def\labelenumi{(\alph{enumi})}
\setcounter{enumi}{1}
\item
  \(A\{S^{-1}\}\) 约化为 0 的充要条件是：在 A 中，0 属于 S 的闭包。举出一个 A 是 Hausdorff 且完备、但 \(S^{-1}A\) 关于 (a) 中定义的拓扑并非如此的例子。
\item
  对于每个连续同态 \(u\)，它从 \(A\) 到完备 Hausdorff 的带线性拓扑环 \(B\)，若 \(u(S)\) 由 \(B\) 的可逆元组成，证明 \(u = u' \circ j\)，其中 \(j: A \to A\{S^{-1}\}\) 是典范映射且 \(u'\) 连续；此外 \(u'\) 唯一确定。
\item
  设 \(S_{1}\)、\(S_{2}\) 是 A 的两个乘性子集，\(S_{2}^{\prime}\) 是 \(S_{2}\) 在 \(A\{S^{-1}\}\) 中的典范像；定义 \(A\{(S_{1}S_{2})^{-1}\}\) 到 \(A\{S_{1}^{-1}\}\{S_{2}'^{-1}\}\) 的典范同构。
\item
  设 \(\mathfrak{a}\) 为 A 的开理想，\(\mathfrak{a}\{\mathrm{S}^{-1}\}\) 为 \(\mathrm{S}^{-1}\mathfrak{a}\) 关于 \(\mathrm{S}^{-1}\mathrm{A}\) 上诱导拓扑的 Hausdorff 完备化；证明 \(\mathfrak{a}\{\mathrm{S}^{-1}\}\) 被典范地识别为 \(\mathrm{A}\{\mathrm{S}^{-1}\}\) 的一个开理想，且离散环 \(\mathrm{A}\{\mathrm{S}^{-1}\} / \mathfrak{a}\{\mathrm{S}^{-1}\}\) 同构于 \(\mathrm{S}^{-1}(\mathrm{A} / \mathfrak{a})\)。反之，若 \(\mathfrak{a}'\) 是 \(\mathrm{A}\{\mathrm{S}^{-1}\}\) 的开理想，则其在 A 中的逆像 \(\mathfrak{a}\) 是一个开理想，且 \(\mathfrak{a}' = \mathfrak{a}\{\mathrm{S}^{-1}\}\)。特别地，映射 \(\mathfrak{p} \mapsto \mathfrak{p}\{\mathrm{S}^{-1}\}\) 将 A 中不与 S 相交的开素理想集递增双射到 \(\mathrm{A}\{\mathrm{S}^{-1}\}\) 的开素理想集。
\item
  设 \(p\) 为 \(A\) 的开素理想，令 \(S = A - p\)。证明 \(A\{S^{-1}\}\) 是局部环，其剩余域同构于 \(A / p\) 的分式域。
\item
  对于每个元素 \(f \in \mathbf{A}\)，令 \(\mathbf{A}_{(f)}\) 表示环 \(\mathbf{A}\{\mathbf{S}_f^{-1}\}\)，其中 \(\mathbf{S}_f\) 是由 \(f^n\)（\(n \geqslant 0\)）组成的乘性集。若 \(f\) 遍历乘性子集 \(\mathbf{S}\)（属于 \(\mathbf{A}\)），则 \(\mathbf{A}_{(f)}\) 构成一个环的直接系统，其直接极限记为 \(\mathbf{A}_{(\mathbf{S})}\)（不带任何拓扑）；定义一个典范同态
\end{enumerate}

\[
\mathrm{A} _ {\{\mathrm{S} \}} \rightarrow \mathrm{A} \{\mathrm{S} ^ {- 1} \}.
\]

 若 \(S = A - p\)，其中 \(p\) 是 \(A\) 的开素理想，证明 \(A_{(S)}\) 是局部环，典范同态 \(A_{(S)} \to A\{S^{-1}\}\) 是局部同态，且 \(A_{(S)}\) 与 \(A\{S^{-1}\}\) 的剩余域典范同构（参见第二章，第 3 节，习题 16）。~

\begin{enumerate}
\def\labelenumi{(\alph{enumi})}
\setcounter{enumi}{7}
\tightlist
\item
   假设 A 上的拓扑是某个 A 的理想 m 的 m-进拓扑，A 是 Hausdorff 且完备的，且 \(m/m^{2}\) 是有限生成的 (A/m)-模。证明若 \(m' = m\{S^{-1}\}\)，则 \(A' = A\{S^{-1}\}\) 上的拓扑是 \(m'\)-进拓扑，\(m' = mA'\)，且 \(m'/m'^{2}\) 是有限生成的 \(A'/m'\)-模（使用第 10 节命题 14）。若 A 是 Noether 环，则 \(A'\) 也是 Noether 环。
\end{enumerate}

\begin{enumerate}
\def\labelenumi{\arabic{enumi}.}
\setcounter{enumi}{27}
\tightlist
\item
  设 A 为带线性拓扑的交换环，E、F 为带线性拓扑的两个拓扑 A-模。若 V（分别 W）遍历 E（分别 F）的开子模集合，则子模
\end{enumerate}

\[
\operatorname{Im} (\mathrm{V} \otimes_ {\mathbf {A}} \mathrm{F}) + \operatorname{Im} (\mathrm{E} \otimes_ {\mathbf {A}} \mathrm{W})
\]

在 \(\mathbf{E} \otimes_{\mathbf{A}} \mathbf{F}\) 中构成某个 0 的邻域基本系，该 \(\mathbf{E} \otimes_{\mathbf{A}} \mathbf{F}\) 的拓扑与其在拓扑环 A 上的模结构相容，称为 E 和 F 上给定拓扑的张量积。此 A-模的 Hausdorff 完备化 \((\mathbf{E} \otimes_{\mathbf{A}} \mathbf{F})^{\wedge}\) 是一个 \(\hat{\mathbf{A}}\)-模，称为 E 与 F 的完备张量积。

\begin{enumerate}
\def\labelenumi{(\alph{enumi})}
\tightlist
\item
  证明若 \((\mathrm{V}_{\lambda})\)（分别 \((\mathrm{W}_{\mu})\)）是 E（分别 F）中由子模组成的 0 的邻域基本系，则 \((\mathbf{E} \otimes_{\mathbf{A}} \mathbf{F})^{\wedge}\) 典范同构于如下 A-模逆系统的逆极限
\end{enumerate}

\[
\left(\mathrm{E} / \mathrm{V} _ {\lambda}\right) \otimes_ {\mathrm{A}} \left(\mathrm{F} / \mathrm{W} _ {\mu}\right);
\]

推出 \((\mathbf{E} \otimes_{\mathbf{A}} \mathbf{F})^{\wedge}\) 是一个 \(\hat{\mathbf{A}}\)-模，且典范同构于 \((\hat{\mathbf{E}} \otimes_{\hat{\mathbf{A}}} \hat{\mathbf{F}})\)；它也记作 \(\hat{\mathbf{E}} \widehat{\otimes}_{\hat{\mathbf{A}}} \hat{\mathbf{F}}\)。

\begin{enumerate}
\def\labelenumi{(\alph{enumi})}
\setcounter{enumi}{1}
\item
  设 \(\mathbf{E}'\)、\(\mathbf{F}'\) 为两个带线性拓扑的拓扑 A-模，且 \(u: \mathbf{E} \to \mathbf{E}', v: \mathbf{F} \to \mathbf{F}'\) 为两个连续 A-线性映射；证明 \(u \otimes v: \mathbf{E} \otimes \mathbf{F} \to \mathbf{E}' \otimes \mathbf{F}'\) 关于 \(\mathbf{E} \otimes \mathbf{F}\) 和 \(\mathbf{E}' \otimes \mathbf{F}'\) 上的张量积拓扑是连续的；将连续线性映射记为 \(u \hat{\otimes} v\)，其为 \((\mathbf{E} \otimes \mathbf{F})^{\wedge} \to (\mathbf{E}' \otimes \mathbf{F}')^{\wedge}\) 且对应于 \(u \otimes v\)。
\item
  设 B、C 为带线性拓扑的两个交换 A-代数，使得典范映射 \(\mathbf{A} \to \mathbf{B}\)、\(\mathbf{A} \to \mathbf{C}\) 连续（为简便起见，称 B、C 为拓扑 A-代数）。证明 \((\mathbf{B} \otimes_{\mathbf{A}} \mathbf{C})^{\wedge}\) 具有典范的拓扑 A-代数结构，称为代数 B、C 的完备张量积。定义典范连续表示
\end{enumerate}

\[
\rho \colon \mathbf {B} \rightarrow (\mathbf {B} \otimes_ {\mathbf {A}} \mathbf {C}) ^ {\wedge}, \quad \sigma \colon \mathbf {C} \rightarrow (\mathbf {B} \otimes_ {\mathbf {A}} \mathbf {C}) ^ {\wedge}
\]

具有如下性质：对于每个完备 Hausdorff 交换拓扑 A-代数 D 以及每一对连续 A-同态 \(u\colon B\to D\)、\(v\colon C\to D\)，存在唯一的连续 A-同态 \(w\colon(B\otimes_{A}C)^{\wedge}\to D\)，使得 \(u=w\circ\rho\) 且 \(v=w\circ\sigma\)。

\begin{enumerate}
\def\labelenumi{(\alph{enumi})}
\setcounter{enumi}{3}
\tightlist
\item
  设 m 为 A 的理想，且 A 上的拓扑是 m-进拓扑。若 E 和 F 都赋予 m-进拓扑，证明在 \(E \otimes_{A} F\) 上，E 和 F 上拓扑的张量积就是 m-进拓扑。推出 \((\mathbf{E} \otimes_{\mathbf{A}} \mathbf{F})^{\wedge}\) 同构于逆极限 \(\varprojlim ((\mathrm{E} / \mathfrak{m}^{n+1}) \otimes_{\mathbf{A}} \mathbf{F})\) 和 \(\varprojlim (\mathbf{E} \otimes_{\mathbf{A}} (\mathbf{F} / \mathfrak{m}^{n+1}\mathbf{F}))\)。
\end{enumerate}

\begin{enumerate}
\def\labelenumi{\arabic{enumi}.}
\setcounter{enumi}{28}
\item
  设 A 为交换环，m 为 A 的理想，且 \(\bigcap_{n\geq0}m^{n}=0\)，c 为 A 中不是零因子的元素。证明传递子 \(q_{n}=m^{n}:Ac\) 关于 m-进拓扑是开的，且其交为 0。若 A 关于 m-进拓扑严格线性紧（习题 22(b)），则 \((q_{n})\) 是此拓扑下 A 中 0 的邻域基本系。
\item
  设 K 为无限交换域，A 为二元不定元的形式幂级数环 \(K[[X,Y]]\)，它是完备的 Hausdorff Noether 局部环。定义 A 的一个乘性子集 S，使得 \(S^{-1}A\) 不是半局部环（若 \((\lambda_{n})\) 是 K 中互异元素的无限序列，考察 A 的素理想 \(p_{n}=A(X+\lambda_{n}Y)\)）。
\item
  证明若 A 是半局部环，则形式幂级数环 A{[}{[}X{]}{]} 也是半局部环（考察此环的 Jacobson 根）。
\end{enumerate}

\exercisesection{3}

\begin{enumerate}
\def\labelenumi{\arabic{enumi}.}
\tightlist
\item
  \begin{enumerate}
  \def\labelenumii{(\alph{enumii})}
  \tightlist
  \item
    设 K 为特征 p \textgreater{} 0 的交换域，B 为如下形式幂级数环：其系数在 K 中，涉及两组无限不定元系 \((\mathbf{X}_{n})\)、\((\mathbf{Y}_{n})\)，且每个给定次数只有有限项（\(\S 2\)，第 12 节）；B 关于 m-进拓扑是 Hausdorff 局部环，m 为其极大理想。令 b 为 B 中由当 \(Y_{i}X_{j}\)、\(i \geqslant 0\) 时的单项式 \(0 \leqslant j \leqslant i\) 生成的闭理想；令 A 为局部环 B/b，它关于 n-进拓扑是 Hausdorff，其中 n = m/b 是其极大理想。令 c 为 A 中等于 \(\sum_{n=1}^{\infty} X_{n}^{p^{n}}\) 模 b 的类的元素；证明当 \(k \geqslant 3\) 时，\(n^{k} \cap (Ac) \not\subset n^{2}c\)（从而关系
  \end{enumerate}
\end{enumerate}

\[
\mathfrak {n} ^ {2} (\mathfrak {n} ^ {2 k} \cap (\mathrm{A} c)) = \mathfrak {n} ^ {2 k + 2} \cap (\mathrm{A} c)
\]

不能成立）。

\begin{enumerate}
\def\labelenumi{(\alph{enumi})}
\setcounter{enumi}{1}
\tightlist
\item
  设 \(C'\) 为底集为 \(A \times A\) 的交换环，其乘法为 \((a, x)(a', x') = (aa', ax' + a'x)\)；令 C 为 \(A \times (Ac)\) 的子环 \(C'\)；C 和 \(C'\) 都是局部环，其极大理想分别记为 r 和 \(r'\)，且 \(C'\) 是有限生成的 C-模。证明 \(r = r' \cap C\)，但 C 上由 \(r'\)-进拓扑诱导的拓扑不是 r-进拓扑。
\end{enumerate}

\begin{enumerate}
\def\labelenumi{\arabic{enumi}.}
\setcounter{enumi}{1}
\tightlist
\item
  设 A 为非域的 Noether 整环，K 为其分式域，m 为 A 的满足 \(\neq\) A 的理想。A 上的 m-进拓扑不是由 K 上的 m-进拓扑诱导的，且后者不是 Hausdorff 的。
\end{enumerate}

* 3. 若 A 是高度 1 的非离散赋值的环（第 VI 章），且 m 是其极大理想，则 A 关于 m-进拓扑不是 Hausdorff 的，\{0\} 的闭包为 m。*

¶ 4. 设 A 为半局部环，r 为其 Jacobson 根。A 为 Noether 环的充要条件是 A 的每个理想关于 r-进拓扑都是闭的，且 A 的每个极大理想都是有限生成的。（若这些条件成立，先使用第 2 节，第 13 节，命题 19 的推论以及第 10 节，定理 1 的推论 5，证明 A 的 Hausdorff 完备化 \(\hat{A}\) 是 Noether 环。然后注意到 r-进拓扑在 A 上是 Hausdorff 的，且 A 的理想集到 \(\hat{A}\) 的理想集存在一个保序单射。）

¶ 5. 设 A 为交换 Noether 环，m 为其 Jacobson 根。

\begin{enumerate}
\def\labelenumi{(\alph{enumi})}
\item
  A 关于线性拓扑 \(\mathcal{T}\) 严格线性紧的充要条件是 A 为半局部环且关于 m-进拓扑完备（\(\S\) 2，习题 19）；此时 \(\mathcal{T}\) 必然与后一个拓扑相同。（使用 \(\S\) 2，习题 21 (d) 和 22 (a)，以及 \(\S\) 3，第 3 节，命题 6。）
\item
  A 关于线性拓扑 T 线性紧的充要条件是 A 为半局部环且关于 m-进拓扑完备，并且 T 比后一个拓扑细。（为看出必要性，先使用第 2 节习题 21 (d) 将其化约到 A 为局部环的情形。先考虑 T 为离散拓扑的情形，注意此时 A 关于 m-进拓扑线性紧。在一般情形下，将其化约到 T 为极小拓扑的情形（第 2 节，习题 18 (c)）；证明此时 A 严格线性紧，并使用第 2 节习题 19 (a) 的判据 (γ) 以及第 2 节习题 18 (b)。）
\end{enumerate}

\begin{enumerate}
\def\labelenumi{\arabic{enumi}.}
\setcounter{enumi}{5}
\item
  设 A 为 Noether 局部整环，m 为其极大理想，K 为其分式域；假设 A 关于 m-进拓扑完备。在 A 上考虑一个比 m-进拓扑细的线性拓扑 T。证明若 \(T'\) 是 K 上的线性拓扑，其 0 的邻域基本系由 A 中关于拓扑 T 的 0 的邻域组成，则 \(T'\) 与 K 的域结构相容，且 K 关于拓扑 \(T'\) 完备。
\item
  设 A 为交换 Noether 环，E 为 A-模。赋予 A 拓扑 \(\mathcal{T}_{m}(\mathrm{A})\)（§ 2，习题 24），并令 \(\mathcal{T}_{m}(\mathrm{E})\) 表示 E 上如下线性拓扑：其 0 的邻域基本系由子模 a.E 组成，其中 a 遍历 \(\mathcal{T}_{m}(\mathrm{A})\) 的 0 的邻域基本系，且这些邻域是 A 的理想。证明若 E 有限生成，则 E 是 Hausdorff 拓扑 A-模（带拓扑 \(\mathcal{T}_{m}(\mathrm{A})\) 和 \(\mathcal{T}_{m}(\mathrm{E})\)），其中每个子模都是闭的。\(\mathcal{T}_{m}(\mathrm{E})\) 的 0 的邻域基本系于是由 E 中使 E/F 成为有限长度模的子模 F 组成，并且对 E 的每个子模 M，拓扑 \(\mathcal{T}_{m}(\mathbf{M})\) 由 \(\mathcal{T}_{m}(\mathbf{E})\) 诱导。
\end{enumerate}

¶ 8. 设 A 为 Noether 半局部环，n 为其幂零根（即 A 的最大幂零理想），m 为其 Jacobson 根。证明若 A（带 m-进拓扑）满足 A/n 完备，则 A 完备。（将其化约到 n 由满足 \(c^{2} = 0\) 的单个元素 c 生成的情形；使用一般拓扑，第 III 章，第 3 节，习题 9，将其化约到证明 n = Ac 完备，并使用 n 是 (A/n)-模这一事实。）

¶ 9. (a) 设 A 为 Zariski 环，m 为 A 的拓扑的定义理想，B 为满足 \(A \subset B \subset \hat{A}\) 且是有限生成 A-模的环；证明若 mB 关于 \(\hat{A}\) 上的拓扑在 B 中是开的，则必有 B = A（使用 Nakayama 引理）。

* (b) 设 A 为交换 Noether 环，m 为 A 的理想。证明若 A 关于 m-进拓扑的 Hausdorff 完备化 \(\hat{\mathbf{A}}\) 是有限生成 A-模，则 A 关于 m-进拓扑完备（将其化约到 A 为 Hausdorff 的情形；使用第五章，第 2 节，第 1 节，命题 1 和定理 1，证明 A 带 m-进拓扑是 Zariski 环，再借助 (a) 得出结论）。*

\begin{enumerate}
\def\labelenumi{\arabic{enumi}.}
\setcounter{enumi}{9}
\item
  设 A 为 Zariski 环，m 为 A 的拓扑的定义理想，E 为有限生成 A-模。证明若 \(\hat{E}\) 是一个 \(\hat{A}\)-模并且有含 r 个元素的生成系，则 A-模 E 有含 r 个元素的生成系（注意 E/mE 与 \(\hat{E}/m\hat{E}\) 同构，并使用第二章，第 3 节，命题 4 的推论 2）。若不假设 E 有限生成，结论是否成立（参见习题 9）？若只假设 A 为 Noether 环，结论是否成立（取 A = Z）？~
\item
  设 A 为 \(\mathbf{Z}_p\)-进整数环 \(p\)（\(p\) 为素数），带使其成为完备 Zariski 环的 \(p\)-进拓扑。令 E 为带 \(\mathrm{A}^{(N)}\)-进拓扑的 A-模 \(p\)。
\end{enumerate}

\begin{enumerate}
\def\labelenumi{(\alph{enumi})}
\item
  证明 \(\hat{\mathbf{E}}\) 的完备化 \(\mathbf{E}\) 被识别为 \(\mathbf{A}^{\mathbf{N}}\) 的子模，该子模由满足 \((a_{n})_{n\in \mathbf{N}}\) 的 \(\mathbf{A}\) 中元素序列 \(\lim a_n = 0\) 组成。
\item
   令 \(e_n\) 为第 \(n\) 个向量，它属于 \(E\) 的典范基。在 \(E\) 中考虑子模 \(F\)，它由向量 \(e_{2n-1}\) 生成，以及子模 \(G\)，它由向量 \(p^{2n}e_{2n} - e_{2n-1}\) 生成（\(n \geqslant 1\)）。证明 \(F\) 和 \(G\) 上的拓扑由 \(p\)-进拓扑在 \(E\) 上诱导，且由 \(p\)-进拓扑诱导的拓扑分别是 \(F\) 和 \(G\) 上的拓扑，但在 \(\hat{E}\) 中，
\end{enumerate}

\[
\overline {{{\mathbf {F} + \mathbf {G}}}} \neq \overline {{{\mathbf {F}}}} + \overline {{{\mathbf {G}}}}.
\]

\begin{enumerate}
\def\labelenumi{(\alph{enumi})}
\setcounter{enumi}{2}
\item
   在 \(\hat{\mathbf{E}}\) 中，令 \(a_{r} = \sum_{n=0}^{\infty} p^{n} e_{r+2}\)（\(r \geqslant 0\)）；令 \(\mathbf{H}\) 为 \(\hat{\mathbf{E}}\) 中由 \(a_{r}\)（\(r \geqslant 0\)）生成的子模。证明由 \(p\)-进拓扑在 \(\hat{\mathbf{E}}\) 上诱导到 H 的拓扑就是 H 上的 \(p\)-进拓扑，但 \(\overline{\mathbf{E} \cap \mathbf{H}} \neq \bar{\mathbf{E}} \cap \bar{\mathbf{H}}\)。
\item
  令 L 为 E 中由 \(p^{n}e_{n}\) 生成的子模；证明 p-进拓扑在 L 上诱导的拓扑不同于 L 上的 p-进拓扑。
\end{enumerate}

\begin{enumerate}
\def\labelenumi{\arabic{enumi}.}
\setcounter{enumi}{11}
\tightlist
\item
  \begin{enumerate}
  \def\labelenumii{(\alph{enumii})}
  \tightlist
  \item
    设 A 为交换 Noether 环，\(m_{1}\)、\(m_{2}\) 为 A 中包含于 A 的 Jacobson 根的两个理想；令 \(A_{i}\) 为 A 关于 \(m_{i}\)-进拓扑的完备化。若 \(m_{1} \subset m_{2}\)，证明 A 上的恒等映射可按连续性延拓为从 \(A_{1}\) 到 \(A_{2}\) 的单射表示（参见一般拓扑，第 III 章，第 3 节，第 5 节，命题 9）。~
  \end{enumerate}
\end{enumerate}

\begin{enumerate}
\def\labelenumi{(\alph{enumi})}
\setcounter{enumi}{1}
\tightlist
\item
  令 \(n = A_1 m_2\)。证明若 \(A_1\) 被典范地识别为 \(A_2\) 的子环，则 \(A_2\) 被识别为 \(A_1\) 关于 \(n\)-进拓扑的完备化。
\end{enumerate}

\begin{enumerate}
\def\labelenumi{\arabic{enumi}.}
\setcounter{enumi}{12}
\tightlist
\item
  设 A 为极大理想为 m 的 Noether 局部环，B 为满足 \(A \subset B \subset \hat{A}\) 的环。假设 B 是极大理想为 n 的 Noether 局部环（习题 14）。
\end{enumerate}

\begin{enumerate}
\def\labelenumi{(\alph{enumi})}
\item
  证明 \(\mathfrak{n} = \mathbf{B} \cap \hat{\mathfrak{m}}\) 且 \(\mathbf{B} = \mathbf{A} + \mathfrak{n}\)。若进一步 \(\mathfrak{n}^2 = \mathbf{B} \cap \hat{\mathfrak{m}}^2\)，则 \(\mathfrak{n} = \mathbf{B}\mathfrak{m}\)（注意此时 \(\mathfrak{n} = \mathbf{B}\mathfrak{m} + \mathfrak{n}^2\)）。
\item
  设 K 为交换域，C 为多项式环 K{[}X{]}，p 为素理想 CX，A 为局部环 \(C_{p}\)（它是 Noether 环），m 为其极大理想；完备化 \(\hat{A}\) 被识别为形式幂级数环 K{[}{[}X{]}{]}（第 4 节，命题 8）。设 \(u = u(X)\) 为 \(\hat{A}\) 中关于有理函数域 K(X) 超越的元素（参见~代数，第 V 章，第 5 节，习题 13，或实变函数，第 III 章，第 1 节，习题 14 (a)），且无常数项；令 B 为 \(\hat{A}\) 的由下列商组成的子环
\end{enumerate}

\[
\mathrm{P} (\mathrm{X}, u (\mathrm{X})) / \mathrm{Q} (\mathrm{X}, u (\mathrm{X})),
\]

其中 P 和 Q 是 K{[}X, Y{]} 中满足 Q(0, 0) ≠ 0 的多项式。证明 B 是包含 A 的 Noether 局部环，其极大理想 n 由 X 和 u(X) 生成；但 B 上的 n-进拓扑严格粗于 m-进拓扑。

\begin{enumerate}
\def\labelenumi{(\alph{enumi})}
\setcounter{enumi}{2}
\tightlist
\item
  在 (b) 的相同定义下，令 \(\mathbf{B}'\) 为由 \(\mathbf{B}\) 和 \(\mathbf{A}\) 生成的 \(u(\mathbf{X})\) 的子环；证明 \(\mathbf{B}'\) 不是局部环（注意 \(\mathbf{B}' \cap \mathfrak{n}\) 是 \(\mathbf{B}'\) 的极大理想，但 \(\mathbf{B}'\) 中存在既不可逆于 \(\mathbf{B}'\) 又不属于 \(\mathbf{B}' \cap \mathfrak{n}\) 的元素）。
\end{enumerate}

¶ 14. 设 K 为交换域，C 为环 K{[}X, Y{]}，p 为 C 的极大理想 CX + CY，A 为局部环 C \(_{p}\)，m 为其极大理想，\(\hat{A}\) 为 A 的完备化，它被识别为形式幂级数环 K{[}{[}X, Y{]}{]}（第 4 节，命题 8）。令 B 为 \(\hat{A}\) 的子环，由形如 Xf(X, Y) + (P(Y)/Q(Y)) 的形式幂级数组成，其中 \(f \in \hat{A}\)，且 P 和 Q 是 K{[}Y{]} 中满足 Q(0) ≠ 0 的两个多项式。

\begin{enumerate}
\def\labelenumi{(\alph{enumi})}
\item
  证明 \(\mathbf{A} \subset \mathbf{B}\)，\(\mathbf{B}\) 是极大理想 \(\mathfrak{n}\) 等于 \(\mathbf{B}\mathfrak{m}\) 的局部环，且对每个整数 \(\mathfrak{n}^k = \mathbf{B} \cap \hat{\mathfrak{m}}^k\)，\(k \geqslant 1\)；因此 \(\mathbf{B}\) 在 \(\hat{\mathbf{A}}\) 中处处稠密，且 \(\mathbf{B}\)-进拓扑在 \(\hat{\mathfrak{m}}\) 上诱导到 \(\hat{\mathbf{A}}\) 的拓扑是 \(\mathfrak{n}\)-进拓扑。
\item
  证明 B 中由所有形如 \(\mathbf{X}f(\mathbf{Y})\) 的元素生成的理想 b（其中 \(f(\mathbf{Y}) \in \mathbf{K}[[\mathbf{Y}]]\)）不是有限生成的（参见~代数，第 V 章，第 5 节，习题 13），因而 B 不是 Noether 环，尽管 B/n 和 \(\mathrm{gr}(\mathbf{B}) = \mathrm{gr}(\hat{\mathbf{A}})\) 是 Noether 环且理想 n 是有限生成的；\(b = B \cap \hat{A}X\)，b 是主理想（非闭）BX 在 B 中的闭包；最后 B 不是平坦 A-模，且 \(\hat{B} = \hat{A}\) 不是平坦 B-模（使用第一章，第 3 节，第 5 节，命题 9）。
\item
   设 \(f(\mathbf{Y})\) 为 \(\mathbf{K}[[\mathbf{Y}]]\) 中的可逆形式幂级数，但不是 \(\mathbf{K}(\mathbf{Y})\) 中的元素。若 \(c\) 是 \(\mathbf{B}\) 中由 \(\mathbf{X}\) 和 \(\mathbf{X}f(\mathbf{Y})\) 生成的理想 \(c\)，证明 \(n\)-进拓扑在 \(c\) 上严格细于 \(n\)-进拓扑在 \(\mathbf{B}\) 上诱导的拓扑。证明典范映射 \(\hat{\mathbf{B}} \otimes_{\mathbf{B}} c \to \hat{c}\)（其中 \(\hat{c}\) 是 \(c\) 关于 \(n\)-进拓扑的完备化）不是单射（考察 \(f(\mathbf{Y}) \otimes \mathbf{X}\) 和 \(1 \otimes \mathbf{X}f(\mathbf{Y})\) 的像；为证明这两个元素在 \(\hat{\mathbf{B}} \otimes_{\mathbf{B}} c\) 中不同，将此张量积看作张量积 \(\hat{\mathbf{B}} \otimes_{\mathbf{K}(\mathbf{Y})} c\) 的商）。
\item
  设 \(f_{1}(\mathbf{Y}), f_{2}(\mathbf{Y})\)、\(\mathbf{K}[[\mathbf{Y}]]\) 为 \(1, f_{1}\) 中的两个可逆形式幂级数，使得 1、\(f_{2}\) 和 \(\mathbf{K}(\mathbf{Y})\) 在 \(c_{1}, c_{2}\) 上线性无关。令 \(\mathbf{B}\)、\(\mathbf{X}f_{1}(\mathbf{Y})\) 分别为 \(\mathbf{X}f_{2}(\mathbf{Y})\) 中由 \(c_{1}\) 和 \(c_{2}\) 生成的主理想；c＿1 和 c＿2 上的 n-进拓扑分别与由 \(\mathbf{B}\) 上拓扑诱导的拓扑相同，而这些理想在 \(\hat{\mathbf{A}} = \hat{\mathbf{B}}\) 中的闭包都与 \(\hat{\mathbf{A}}\mathbf{X}\) 的主理想 \(\hat{\mathbf{A}}\) 相同。推出 \(c_{1} \cap c_{2}\) 在 \(\hat{\mathbf{A}}\) 中的闭包不等于 \(\bar{c}_{1} \cap \bar{c}_{2}\)，且 \(c_{1}: c_{2}\) 的闭包不等于 \(\bar{c}_{1}: \bar{c}_{2}\)。
\end{enumerate}

\begin{enumerate}
\def\labelenumi{\arabic{enumi}.}
\setcounter{enumi}{14}
\tightlist
\item
  \begin{enumerate}
  \def\labelenumii{(\alph{enumii})}
  \tightlist
  \item
    设 A 为 Noether 整环，m 为 A 的理想，\(\hat{A}\) 为 A 关于 m-进拓扑的 Hausdorff 完备化。证明若 M 是无挠 A-模，则对于 \(\hat{A}\) 中每个不是零因子的元素 b，比为 b 的倍乘在 \(\hat{A}\)-模 \(\hat{A} \otimes_{A} M\) 上是单射。（将其化约到 M 有限生成的情形；此时存在一个极大
  \end{enumerate}
\end{enumerate}

 自由系 \((m_j)_{1\leqslant j\leqslant n}\)，并且存在 \(a\in \mathbf{A}\)，使得对所有 \(m\in \mathbf{M}\)，\(am = \sum_{j}a_{j}m_{j}\)，其中 \(a_{j}\in \mathbf{A}\)；因此每个 \(x\in \hat{\mathbf{A}}\otimes_{\mathbf{A}}\mathbf{M}\) 都满足 \(ax = \sum_{j}b_{j}\otimes m_{j}\)，其中 \(b_{j}\in \hat{\mathbf{A}}\)；然后使用 \(\hat{\mathbf{A}}\) 作为 A-模的平坦性。

\begin{enumerate}
\def\labelenumi{(\alph{enumi})}
\setcounter{enumi}{1}
\tightlist
\item
  设 K 为特征 0 的代数闭域，B 为多项式环 K{[}X, Y{]}，且 P(X, Y) = X(X² + Y²) + (X² - Y²)。证明理想 BP 在 B 中是素理想；考察商环 A = B/BP，它是 Noether 整环。令 m 为 A 的极大理想，即 B 的极大理想 n = BX + BY 的典范像。证明 A 关于 m-进拓扑的 Hausdorff 完备化 \(\hat{A}\) 不是整环（注意在形式幂级数环 K{[}{[}X, Y{]}{]} 中，P 分解为两个形式幂级数的乘积（``不可约三次曲线的双重点''））。
\end{enumerate}

\begin{enumerate}
\def\labelenumi{\arabic{enumi}.}
\setcounter{enumi}{15}
\tightlist
\item
  设 A 为交换 Noether 环，m 为 A 的理想，E 为带 m-进拓扑的有限生成 A-模。对于每个无限集 I，令 \(E_{I}\) 表示 E 中元素族 \((x_{i})_{i\in I}\) 的集合，其中 lim \(x_{i}=0\) 是相对于 I 的有限子集补集滤子而言的；它是乘积 A-模 \(E^{I}\) 的子模。
\end{enumerate}

\begin{enumerate}
\def\labelenumi{(\alph{enumi})}
\item
  证明若 \(0 \rightarrow E' \rightarrow E \rightarrow E'' \rightarrow 0\) 是有限生成 A-模的正合列，则相应的序列 \(0 \rightarrow E_{1}' \rightarrow E_{1} \rightarrow E_{1}'' \rightarrow 0\) 也是正合的。
\item
  对每个 A-模 E 定义一个典范 A-同态 \(A_{I} \otimes_{A} E \to E_{1}\)，并证明它是同构（先在 E 自由且有限生成时验证，再使用 (a)）。
\item
  由 (b) 推出 \(A_{I}\) 是忠实平坦 A-模。
\end{enumerate}

¶ 17. (a) 设 K 为交换域，A = K{[}{[}X₁, \ldots, Xₙ{]}{]} 为系数在 K 中的 n 个不定元的形式幂级数环，V 为 K 上的向量空间。用第 2 节，第 6 节，例 1 的记号，令 V{[}{[}X₁, \ldots, Xₙ{]}{]} 表示向量空间 Vⁿ，并赋予如下定义的 A-模结构：

\[
\left(\sum_ {\alpha} c _ {\alpha} \mathbf {X} ^ {\alpha}\right) (v _ {\alpha}) = (w _ {\alpha}),
\]

其中 \(w_{\alpha} = \sum_{\beta + \gamma = \alpha} c_{\beta} v_{\gamma}\)。证明若 V 有一个以 I 为指标集的基，则 A-模 \(V[[X_1, \ldots, X_n]]\) 在 I 有限时同构于 \(A^I\)，在 I 无限时同构于习题 16 中定义的 A-模 \(A_1\)。推出 \(V[[X_1, \ldots, X_n]]\) 是平坦 A-模，并且当 V 不约化为 0 时是忠实平坦的。

\begin{enumerate}
\def\labelenumi{(\alph{enumi})}
\setcounter{enumi}{1}
\item
  设 L 是 K 的扩域。由 (a) 推出形式幂级数环 \(L[[X_{1},\ldots,X_{n}]]\) 是环 \(K[[X_{1},\ldots,X_{n}]]\) 上的忠实平坦模。若 L 在 K 上秩有限，则 \(L[[X_{1},\ldots,X_{n}]]\) 同构于 \(L\otimes_{K}K[[X_{1},\ldots,X_{n}]]\)。
\item
  若 \(\mathbf{L}\) 是 \(\mathbf{K}\) 的代数扩张，证明形式幂级数环 \(\mathbf{L}[[\mathbf{X}_1, \ldots, \mathbf{X}_n]]\) 是环
\end{enumerate}

\[
\mathbf {L} \otimes_ {\mathbf {K}} \mathbf {K} [ [ \mathbf {X} _ {1}, \dots , \mathbf {X} _ {n} ] ]
\]

（将 L 看作其在 K 上有限秩子扩张的直接极限，并使用第一章，第 2 节，第 7 节，命题 9）。推出环

\[
\mathbf {B} = \mathbf {L} \otimes_ {\mathbf {K}} \mathbf {K} [ [ \mathbf {X} _ {1}, \dots , \mathbf {X} _ {n} ] ]
\]

是一个 Noether 局部环，其完备化被识别为 \(\mathbf{L}[[\mathbf{X}_1,\ldots ,\mathbf{X}_n]]\)。要使 \(\hat{\mathbf{B}} = \mathbf{B}\)，充要条件是 \(n = 0\) 或 \([\mathbf{L}:\mathbf{K}] < +\infty\)。

¶ 18. 设 A 为极大理想为 m 的局部环；若 A-模 M 使 M/mM 是剩余域 \(k = \mathbf{A} / \mathfrak{m}\) 上的有限秩向量空间，则称 M 为拟有限的。特别地，若 \(\mathbf{A}\) 是整环，则 \(\mathbf{K}\) 的分式域 \(\mathbf{A}\) 是拟有限的 \(\mathbf{A}\)-模。

\begin{enumerate}
\def\labelenumi{(\alph{enumi})}
\item
  证明若 A 是 Noether 环且 M 是拟有限 A-模，则 M 关于 m-进拓扑的 Hausdorff 完备化 \(\hat{\mathbf{M}}\) 是有限生成的 \(\hat{\mathbf{A}}\)-模（使用第 2 节，第 11 节，命题 14 的推论 2）。特别地，若 A 完备且 M 关于 m-进拓扑为 Hausdorff，则 M 是有限生成的 A-模。
\item
  设 B 为另一个局部环，n 为其极大理想，\(\phi: A \rightarrow B\) 为局部同态，M 为有限生成 B-模。证明若 B 是 Noether 环且 M 是拟有限 A-模，则 M 上的 m-进拓扑和 n-进拓扑相同。（注意 M/mM 是有限长度 B-模，并推出若 b 是 B-模 M/mM 的湮灭子，则 \(V(b) = \{n\}\) 在 Spec(B) 中成立，因而 \(V(mB + b) = \{n\}\)。使用第二章，第 4 节，第 3 节，命题 11 的推论 2 得出结论。）
\item
  在 (b) 的假设下，证明若 \(M \neq 0\)，B/b 是拟有限 A-模（注意 \(M \neq nM\)，且 M/nM 是 k 上的有限秩向量空间；推出 B/n 在 k 上秩有限）。
\end{enumerate}

¶ 19. 设 A 为交换 Noether 环，m 为 A 的理想，S 为 A 的乘性子集。赋予 A m-进拓扑。

\begin{enumerate}
\def\labelenumi{(\alph{enumi})}
\item
  证明环 \(\mathbf{A}\{\mathbf{S}^{-1}\}\)（第 2 节，习题 27）是平坦 A-模。
\item
   设 \(S'\) 是 \(A\) 的另一个乘性子集，包含于 \(S\)。证明 \(A\{S^{-1}\}\) 是平坦的 \((A\{S'^{-1}\})\)-模（使用第 2 节习题 27 (d)）。
\item
  证明 \(\mathbf{A}\{\mathbf{S}^{-1}\}\) 是平坦的 \(\mathbf{A}_{\{\mathbf{S}\}}\)-模（第 2 节，习题 27 (g)）。
\item
   假设 \(S = A - p\)，其中 \(p\) 是 \(A\) 的开素理想。证明 \(A\{S^{-1}\}\) 是忠实平坦的 \(A_{(S)}\)-模，并推出环 \(A_{(S)}\) 是 Noether 环（参见第 2 节，习题 27 (g)）。~
\end{enumerate}

\begin{enumerate}
\def\labelenumi{\arabic{enumi}.}
\setcounter{enumi}{19}
\item
  设 A 为交换环，m 为 A 的理想；赋予 A m-进拓扑。设 B 为交换拓扑 A-代数（\(\S 2\)，习题 28）；假设 B 是 Zariski 环；令 n 为 B 的定义理想。证明若 M 是带 m-进拓扑的有限生成 A-模，则在 B-模 \(B \otimes_{A} M\) 上，B 上的拓扑与 M 上的 m-进拓扑的张量积就是 n-进拓扑；于是完备张量积 \((\mathbf{B} \otimes_{\mathbf{A}} \mathbf{M})^{\wedge}\) 同构于 \(\hat{B} \otimes_{A} M\)。
\item
  设 A 为交换 Noether 环，m 为 A 的理想，M、N 为两个有限生成 A-模。
\end{enumerate}

\begin{enumerate}
\def\labelenumi{(\alph{enumi})}
\item
  假设 M 关于 m-进拓扑是 Hausdorff 的。证明在带 m-进拓扑的 \(\operatorname{Hom}_{\mathbf{A}}(\mathbf{M}, \mathbf{N})\) 中，单射同态的集合是开的（使用第 1 节，命题 2 和 Artin–Rees 引理）。
\item
  假设 A 是完备 Zariski 环，m 是其拓扑的定义理想
\end{enumerate}

A。对每个整数 \(i\)，令 \(\mathbf{A}_i = \mathbf{A} / \mathfrak{m}^{i+1}\)、\(\mathbf{M}_i = \mathbf{M} / \mathfrak{m}^{i+1}\mathbf{M}\)、\(\mathbf{N}_i = \mathbf{N} / \mathfrak{m}^{i+1}\mathbf{N}\)；证明拓扑 A-模 \(\operatorname{Hom}_{\mathbf{A}}(\mathbf{M}, \mathbf{N})\) 同构于

\[
\varprojlim \operatorname{Hom} _ {\mathbf {A}} (\mathbf {M} _ {i}, \mathbf {N} _ {i}).
\]

推出在 \(\operatorname{Hom}_{\mathbf{A}}(\mathbf{M},\mathbf{N})\) 中，满射同态的集合是开的。

¶ 22. (*) 设 A 为环，无论交换与否；以下所考虑的所有 A-模都是左 A-模。设 P 为一种性质，满足：(α) 若 f: M → N 是单射 A-模同态且 N 具有性质 P，则 M 具有性质 P；(β) 具有性质 P 的两个 A-模的直和仍具有性质 P。

\begin{enumerate}
\def\labelenumi{(\alph{enumi})}
\item
   设 M 为 A-模。证明子模 \(M'\)（满足 \(M/M'\) 具有性质 P）构成 M 上线性拓扑 \(\mathcal{T}_{P}(M)\) 的 0 的邻域基本系。证明 \(\mathcal{T}_{P}(A_{s})\) 与 A 的环结构相容，且带拓扑 \(\mathcal{T}_{P}(M)\) 的 M 是带 \(\mathcal{T}_{P}(A_{s})\) 的 A 上的拓扑模。每个 A-模同态 \(f: M \to N\) 关于拓扑 \(\mathcal{T}_{P}(M)\) 和 \(\mathcal{T}_{P}(N)\) 都是连续的。
\item
  假设满足以下条件：\((\gamma)\) 若 \(N\) 是 \(M\) 的具有性质 \(P\) 的子模，且对于每个子模 \(L \neq 0\)（它是 \(M\) 的子模），都有 \(N \cap L \neq 0\)，则 \(M\) 具有性质 \(P\)。证明在这些条件下，若 \(F\) 是 A-模 \(E\) 的子模，则拓扑 \(\mathcal{T}_P(F)\) 由 \(\mathcal{T}_P(E)\) 诱导。（令 \(F'\) 为 \(F\) 的一个子模，使 \(F/F'\) 具有性质 \(P\)；考察 \(G\)，它是 \(E\) 的子模且满足 \(G \cap F = F'\)，并证明 \(E/G\) 具有性质 \(P\)。）
\end{enumerate}

¶ 23. (a) 设 A 为交换 Noether 环，m 为 A 的理想。令 P\{M\} 表示如下性质：M 是 A-模，且 M 的每个有限生成子模都被 m 的某个幂湮灭。

证明习题 22 的条件 \((\alpha)\)、\((\beta)\) 和 \((\gamma)\) 得到满足。（为证明 \((\gamma)\)，将其化约到 M 有限生成的情形；对于所有 \(a \in \mathfrak{m}\)，由假设存在 \(k > 0\) 使 \(a^k\mathrm{N} = 0\)；使用如下事实：存在 \(r > 0\) 使 \(\operatorname{Ker}(a_{\mathbf{M}}^r) \cap \operatorname{Im}(a_{\mathbf{M}}^r) = 0\)（代数，第 VIII 章，第 2 节，第 2 节，引理 2）。）

\begin{enumerate}
\def\labelenumi{(\alph{enumi})}
\setcounter{enumi}{1}
\item
  证明若 M 是有限生成 A-模，则拓扑 \(\mathcal{T}_{P}(\mathbf{M})\) 与 m-进拓扑相同。给出一个 A-模 M 的例子，使 \(\mathcal{T}_{P}(\mathbf{M})\) 严格细于 m-进拓扑（参见习题 11）。~
\item
  证明 (a) 的结论不能推广到 A 是非交换左 Noether 环且 m 是 A 的双边理想的情形（考察交换域上二阶下三角矩阵环）。
\end{enumerate}

¶ 24. 若环 A（无论交换与否）不约化为 0，且没有零因子、A 的每个左理想或右理想都是单生成的，则称 A 为主理想环。这样的环既是左 Noether 环也是右 Noether 环。

\begin{enumerate}
\def\labelenumi{(\alph{enumi})}
\item
   证明对于每个元素 \(c \neq 0\) 的 \(\mathbf{A}\)，\(\mathbf{A} / \mathbf{A}c\) 是有限长度的 \(\mathbf{A}\)-模。（注意若理想递减序列 \(\mathbf{A}a_n\)（在 \(\mathbf{A}\) 中）包含 \(\mathbf{A}c\)，则有 \(c = b_na_n\) 对所有 \(n\) 成立，并考察右理想 \(b_na\)。）
\item
  证明（左或右）自由 A-模的每个子模都是自由的（论证同代数，第 VII 章，第 3 节，定理 1）。
\item
  在每个 A-模 M 中，M 的非自由元素构成 M 的一个子模 T，称为 M 的挠子模（使用第二章，第 2 节，习题 14 (a)）；若 T = M，则称 M 为挠模；若 T = 0，则称 M 为无挠的。
\item
  证明每个有限生成的无挠 A-模都是自由的（使用第二章，第 2 节，习题 23 (b)）。
\item
  证明每个有限生成 A-模都是一个自由模和一个挠模的直和。
\item
   设 \(\mathfrak{a}\) 为 A 的理想 \(\neq 0\)。证明存在元素 \(a \in \mathfrak{a}\) 以及 A 的自同构 \(\sigma\)，使得 \(\mathfrak{a} = \mathrm{A}a = a\mathrm{A}\)，且满足 \(ax = \sigma(x)a\)，对所有 \(x \in \mathrm{A}\) 都有。（若 \(b\) 是左理想的 \(a\)-生成元，所生成的理想为 \(\mathfrak{a}\)，则存在 A 的自同态 \(\tau\)，使得 \(bx = \tau(x)b\) 对所有 \(x \in \mathrm{A}\) 都成立；使用代数，第 VIII 章，第 2 节，习题 8 (b)，证明若 \(a = ub\) 是右理想 \(\mathfrak{a}, u\) 的一个生成元，则可逆。）
\end{enumerate}

¶ 25. 设 A 为主理想环（习题 24），a 为 A 的非零双边理想，且 a 为 a 中具有习题 24 (f) 所述性质的一个元素。令 P\{M\} 表示如下性质：M 是左 A-模，且 M 的每个有限生成子模都被 a 的某个幂湮灭。

\begin{enumerate}
\def\labelenumi{(\alph{enumi})}
\item
   证明习题 22 的条件 \((\alpha)\)、\((\beta)\) 和 \((\gamma)\) 得到满足（考察倍乘 \(a_{\mathbf{M}}\)，仿照习题 23 (a) 进行论证，并注意 \(\operatorname{Ker}(a_{\mathbf{M}}^{r})\) 和 \(\operatorname{Im}(a_{\mathbf{M}}^{r})\) 是 \(\mathbf{M}\) 的子模）。
\item
  证明每个挠 A-模 M（习题 24 (c)）都是一个具有性质 P 的子模 \(M_{a}\) 与一个子模 \(M_{a}^{\prime}\) 的直和，其中倍乘 \(M_{a}^{\prime}\) 在 \(a_{M}\) 上的限制是双射（注意若 N 是挠 A-模且 \(a_{N}\) 是单射，则 \(a_{N}\) 是双射；为此将其化约到 N 单生成的情形，并使用习题 24 (a) 以及代数，第 VIII 章，第 2 节，第 2 节，引理 2）。
\item
   令 S 为 A 中满足其在环 A/a 中的典范像可逆的元素 \(s \in A\) 的集合。证明 S 是 A 的乘性子集，且对于 A 的任意左理想 l，下列条件等价于：
\end{enumerate}

\((\alpha) \mathfrak{l} \cap \mathbb{S} \neq \varnothing; (\beta) \mathfrak{l} + \mathfrak{a} = \mathbb{A}; (\gamma) (\mathbb{A} / \mathfrak{l})_{\mathfrak{a}} = 0\)（使用 (b) 的记号）；(δ) 对所有 \(x \in \mathbb{A}\)，存在 \(s \in \mathbb{S}\)，使得 \(sx \in \mathfrak{l}\)。

推出 A 关于 S 有一个（左或右）分式环（第二章，第 2 节，习题 22），它是主理想环，且其唯一的非零双边理想由理想 \(a^{n}\) 的典范像生成；此外，A 到此分式环的典范映射是单射。

\begin{enumerate}
\def\labelenumi{(\alph{enumi})}
\setcounter{enumi}{3}
\tightlist
\item
  现在假设双边理想 a 是极大的。证明对于每个整数 n \textgreater{} 0，环 \(A/a^{n}\) 同构于完全准素环上的矩阵环 \(\mathbf{M}_{r}(\mathbf{B}_{n})\)（代数，第 VIII 章，第 6 节，习题 20）。若 \(\mathfrak{b}_n\) 是 \(\mathbf{B}_n\) 的极大理想，证明 \(\mathfrak{b}_n^n = 0\)，且 \(\mathbf{B}_n\) 的每个（左或右）理想都形如 \(\mathfrak{b}_n^k\)，并且是单生成的。（一方面注意
\end{enumerate}

\[
\mathbf {a} / \mathbf {a} ^ {n} = \mathbf {M} _ {r} (\mathfrak {b} _ {n})
\]

（代数，第 VIII 章，第 6 节，习题 5），另一方面，对 \(k < n\)，\(\mathfrak{b}_n^k / \mathfrak{b}_n^{k+1}\) 必然是简单 \(B_n\)-模，否则 \(\mathbf{M}_r(\mathfrak{b}_n^k)\) 不可能是单生成的 \((A / a^n)\)-模。）

 推出 \(\hat{\mathbf{A}}\) 是 \(\mathbf{A}\) 关于拓扑 \(\mathcal{T}_{\mathbf{P}}(\mathbf{A}_s)\) 的完备化，即为矩阵环 \(\mathbf{M}_r(\mathbf{B})\)，其中 \(\mathbf{B}\) 没有零因子，且其每个理想（左理想或右理想）都是同一个极大双边理想的幂。

\begin{enumerate}
\def\labelenumi{\arabic{enumi}.}
\setcounter{enumi}{25}
\tightlist
\item
  设 B 为交换环，A 为 B 的完备 Noether 半局部子环。设 n 为 B 的理想，包含 A 的 Jacobson 根的某个幂，且 B 上的 n-进拓扑是 Hausdorff 的。证明 A 上的拓扑由 B 上的 n-进拓扑诱导（使用第 2 节，第 7 节，命题 8）。
\end{enumerate}

¶ 27. 设 A 为环，B 为交换 A-代数且是 Zariski 环，N 为有限生成 B-模。

\begin{enumerate}
\def\labelenumi{(\alph{enumi})}
\item
  假设 B 的理想 \(\mathfrak{I}\) 包含于 B 的 Jacobson 根，且对 \(\mathbf{N} / \mathfrak{I}^{n + 1}\mathbf{N}\)，A-模 \(n\geqslant 0\) 都是平坦的。证明 N 是平坦 A-模。（若 \(v\colon \mathbf{M}\to \mathbf{M}'\) 是有限生成 A-模的单射同态，则必须证明 \(u = v\otimes 1:\mathbf{M}\otimes_{\mathbf{A}}\mathbf{N}\rightarrow \mathbf{M}'\otimes_{\mathbf{A}}\mathbf{N}\) 是单射。将其化约到证明如下命题：若记 \(\mathbf{N}_n = \mathbf{N} / \mathfrak{I}^{n + 1}\mathbf{N}\) 且 \(u_{n} = v\otimes 1_{\mathbf{N}_{n}}\)，则同态 \(\lim u_n\) 是单射；这里使用 \(\mathfrak{I}\) 和 \(\mathbf{M}\otimes_{\mathbf{A}}\mathbf{N}\) 上的 \(\mathbf{M}'\otimes_{\mathbf{A}}\mathbf{N}\)-进拓扑是 Hausdorff 这一事实。）
\item
   设 \(b\) 为包含于 \(\mathbf{B}\) 的 Jacobson 根中的元素，且乘以 \(b\) 在 \(\mathbf{N}\) 上是单射。证明若 \(\mathbf{N} / b\mathbf{N}\) 是平坦 A-模，则 \(\mathbf{N}\) 是平坦 A-模（将其化约到 (a) 中的情形）。
\item
  进一步假设 A 是极大理想为 m、剩余域为 \(k = A/m\) 的局部环，且 mB 包含于 B 的 Jacobson 根。令 P 为平坦 A-模的 B-模，\(u: N \to P\) 为满足 \(u \otimes 1_k: N \otimes_A k \to P \otimes_A k\) 是单射的同态。证明 N 是平坦 A-模且 u 是单射。（将其化约到证明对于每个有限生成 A-模 M，同态 \(u \otimes 1_M: N \otimes_A M \to P \otimes_A M\) 是单射；注意 \(N \otimes_A M\) 上的 m-进拓扑是 Hausdorff 的；然后使用第 2 节，第 8 节，定理 1 的推论 1，考察交换图
\end{enumerate}

\[
\begin{array}{c} \operatorname{gr} _ {0} (\mathbf {N}) \otimes_ {k} \operatorname{gr} (\mathbf {M}) \xrightarrow {\operatorname{gr} _ {0} (u) \otimes 1} \operatorname{gr} _ {0} (\mathbf {P}) \otimes_ {k} \operatorname{gr} (\mathbf {M}) \\ \Biggl \downarrow \\ \operatorname{gr} (\mathbf {N} \otimes_ {\mathbf {A}} \mathbf {M}) \xrightarrow [ \operatorname{gr} (u \otimes 1) ]{} \operatorname{gr} (\mathbf {P} \otimes_ {\mathbf {A}} \mathbf {M}) \end{array}
\]

并使用 P 的平坦性，其中竖直箭头是第 2 节习题 6 的典范同态。）

\exercisesection{4}

\begin{enumerate}
\def\labelenumi{\arabic{enumi}.}
\tightlist
\item
  设 A 为交换环，带滤过 \((a_{n})_{n\geq0}\) 且关于此滤过是 Hausdorff 且完备的，满足 \(a_{0}=A\)。令 M、\(M'\)、N 为三个 A-模，每个都赋予由 A 上滤过诱导的滤过以及由该滤过定义的拓扑；记 \(\overline{M}=M/a_{1}M\)、\(\overline{M}'=M'/a_{1}M'\)、\(\overline{N}=N/a_{1}N\)。令 \(f: M\times M'\to N\) 为双线性映射，\(\bar{f}: \overline{M}\times\overline{M'}\to\overline{N}\) 为由 f 取商得到的 \((A/a_{1})\)-双线性映射。令 \(y\in N\)、\(\bar{x}\in\overline{M}\)、\(\bar{x}'\in\overline{M'}\) 满足：(1) \(\bar{f}(\bar{x},\bar{x}')\) 是 y 在 \(\bar{y}\) 中的类 \(\overline{N}\)；(2) \(\overline{N}\) 的每个元素都可写成
\end{enumerate}

\[
\bar {f} (\overline {{{x}}}, \overline {{{z}}} ^ {\prime}) + \bar {f} (\overline {{{z}}}, \overline {{{x}}} ^ {\prime}),
\]

其中 \(\bar{z} \in \overline{\mathbf{M}}\)，\(\bar{z}' \in \overline{\mathbf{M}}'\)。证明若 N 是 Hausdorff 的且 M、M' 完备，则存在 \(x \in \bar{x}\) 和 \(x' \in \bar{x}'\)，使得 \(f(x, x') = y\)（仿照 Hensel 引理的证明作归纳）。在什么条件下 \(x\) 和 \(x'\) 被唯一确定？

\begin{enumerate}
\def\labelenumi{\arabic{enumi}.}
\setcounter{enumi}{1}
\tightlist
\item
  设 A 为局部环，m 为其极大理想，k = A/m 为其剩余域，\(f: A \to k\) 为典范同态。令 \(P \in A[X]\) 为次数为 n 的首一多项式。~记 \(B = A[X]/P \cdot A[X]\)，并以 x 表示 X 在 B 中的典范像。
\end{enumerate}

\begin{enumerate}
\def\labelenumi{(\alph{enumi})}
\item
   设 \(\mathbf{Q}, \mathbf{Q}'\) 是 \(\mathbf{A}[\mathbf{X}]\) 中两个强互素的首一多项式，满足 \(\mathbf{P} = \mathbf{QQ}'\)。证明环 \(\mathbf{B}\) 是理想 \(\mathbf{B}. \mathbf{Q}(x)\) 和 \(\mathbf{B}. \mathbf{Q}'(x)\) 的直和。
\item
  反之，设 \(B = b \oplus b'\) 是 B 作为两个理想直和的分解。证明存在多项式 Q、\(Q'\) 在 A{[}X{]} 中，满足 (a) 的假设，并且 \(\mathfrak{b} = \mathrm{B}.Q(x)\)、\(\mathfrak{b}' = \mathrm{B}.Q'(x)\)。（先证明 b/mb 和 \(b'/mb'\) 由首一多项式 \(Q_{0} \in \bar{f}^{-1}(\mathfrak{b})\)、\(Q'_{0} \in \bar{f}^{-1}(\mathfrak{b}')\) 在 B/mB 中的像生成，并满足
\end{enumerate}

\[
\bar {f} (\mathrm{P}) = \bar {f} (\mathrm{Q} _ {0}) \bar {f} (\mathrm{Q} _ {0} ^ {\prime}).
\]

令 \(r = \deg(Q_0)\)、\(s = \deg(Q_0')\)。证明 b 是以 \(Q_0(x)\)、\(xQ_0(x)\)、\(\ldots\)、\(x^{s-1}Q_0(x)\) 为基的自由 A-模（使用第二章，第 3 节，第 2 节，命题 5）；于是可写成 \(x^sQ_0(x) = a_0Q_0(x) + a_1xQ_0(x) + \cdots + a_{s-1}x^{s-1}Q_0(x)\)，其中 \(a_i \in A\)（\(0 \leqslant i \leqslant s-1\)）；证明若写成

\[
\mathrm{Q} ^ {\prime} (\mathrm{X}) = \mathrm{X} ^ {s} - (a _ {0} + a _ {1} \mathrm{X} + \dots + a _ {s - 1} \mathrm{X} ^ {s - 1}),
\]

则 \(\bar{f}(\mathbf{P}) = \bar{f}(\mathbf{Q}_0)\bar{f}(\mathbf{Q}')\)，且 \(\mathbf{Q}'\in \bar{f}^{-1}(\mathfrak{b}')\)。同样从 \(\mathbf{Q}\) 和 \(\mathbf{Q}'\) 出发定义 \(\mathbf{Q}_0\)，并使用第 1 节命题 2，证明 \(\mathbf{Q}\) 和 \(\mathbf{Q}'\) 解决此问题。）

¶ 3. 设 A 为局部环，m 为其极大理想，k = A/m 为其剩余域，f: A → k 为典范同态。证明以下两个条件等价：

\begin{enumerate}
\def\labelenumi{(\Alph{enumi})}
\setcounter{enumi}{7}
\item[(H)]
  对每个首一多项式 \(\mathbf{P} \in \mathbf{A}[\mathbf{X}]\)，以及 \(\bar{f}(\mathbf{P}) \in k[\mathbf{X}]\) 作为互素首一多项式之积 \(\bar{f}(\mathbf{P}) = \overline{\mathbf{Q}}\)。\(\overline{\mathbf{Q}}'\) 的每个分解，存在 \(\mathbf{Q}\) 中两个首一多项式 \(\mathbf{Q}'\)、\(\mathbf{A}[\mathbf{X}]\)，使得 \(\bar{f}(\mathbf{Q}) = \overline{\mathbf{Q}}\)、\(\bar{f}(\mathbf{Q}') = \overline{\mathbf{Q}}'\)，且 \(\mathbf{P} = \mathbf{QQ'}\)。
\item[(C)]
  每个作为有限生成 A-模的 A 上交换代数，都是局部 A-代数的直积。
\end{enumerate}

（为证明 (H) 蕴含 (C)，仿照第 6 节命题 8，使用习题 2 (a)。为看出 (C) 蕴含 (H)，先证明对于每个作为有限生成 A-模的交换 A-代数 B，B/mB 作为两个理想的直和的每个分解必然形如 b/mb ⊕ b'/mb'，其中 B = b ⊕ b' 是 B 作为两个理想的直和分解；然后使用习题 2 (b)。）

满足条件 (H) 和 (C) 的局部环称为 Hensel 环。每个完备 Hausdorff 局部环都是 Hensel 环。若 A 是 Hensel 环，B 是交换 A-代数、局部环且为有限生成 A-模，则 B 是 Hensel 环。

\begin{enumerate}
\def\labelenumi{\arabic{enumi}.}
\setcounter{enumi}{3}
\tightlist
\item
  \begin{enumerate}
  \def\labelenumii{(\alph{enumii})}
  \tightlist
  \item
    设 \((\mathbf{A}_{\alpha}, \phi_{\alpha \beta})\) 为 Hensel 局部环的直接系统，同态 \(\phi_{\alpha \beta}\) 为局部同态。证明局部环 \(\mathbf{A} = \varinjlim \mathbf{A}_{\alpha}\)（第二章，第 3 节，习题 16）是 Hensel 环（使用习题 3 的判据 (H)）。
  \end{enumerate}
\end{enumerate}

\begin{enumerate}
\def\labelenumi{(\alph{enumi})}
\setcounter{enumi}{1}
\tightlist
\item
  设 K 为交换域，L 为 K 的代数扩张。由 (a) 推出环 \(L \otimes_{K} K[[X_{1}, \ldots, X_{n}]]\) 是 Hensel 环。* (c) 设 A 为 Hensel 局部环，B 为在 A 上整的交换 A-代数（第五章），且 B 是局部环。证明 B 是 Hensel 环（使用 (a)）。*
\end{enumerate}

¶ 5. 设 A 为 Hensel 局部环，B 为有限生成 A-模的（不必交换的）A-代数，b 为 B 的双边理想，令 \(\overline{\mathbf{B}} = \mathbf{B} / \mathfrak{b}\)。

\begin{enumerate}
\def\labelenumi{(\alph{enumi})}
\item
  证明 \(\varepsilon\) 中的每个幂等元 \(\overline{\mathbf{B}}\) 都是 \(\mathbf{B}\) 中某个幂等元的典范像（考虑由单个元素生成的 \(\mathbf{B}\) 的子代数，将其化约到交换情形）。
\item
  令 \((\varepsilon_{n})_{n\geqslant 1}\) 为 \(\overline{\mathbf{B}}\) 中元素的无限序列，使得
\end{enumerate}

\[
\varepsilon_ {i} \varepsilon_ {j} = \delta_ {i j} \varepsilon_ {j}
\]

对于每个指标有序对 \((i,j)\)。证明 B 中存在幂等元的正交序列 \((e_n)_{n\geq 1}\)，使得 \(\varepsilon_{n}\) 都是 \(e_n\) 的典范像，对所有 \(n\)。（仿照代数，第 VIII 章，第 6 节，习题 10 作归纳；注意若 \(e,e^{\prime}\) 是 B 的两个满足 \(ee^{\prime} = 0\) 的幂等元，则 \(e^{\prime} - e^{\prime}e = e''\) 是满足 \(ee'' = e''e = 0\) 的幂等元。）

\begin{enumerate}
\def\labelenumi{(\alph{enumi})}
\setcounter{enumi}{2}
\tightlist
\item
  现在假设 b 是 B 的 Jacobson 根。令 n 为一个整数，\((\varepsilon_{ij})\)（\(1 \leqslant i \leqslant n, 1 \leqslant j \leqslant n\)）为 \(\overline{B}\) 中的元素族，满足 \(\varepsilon_{ij}\varepsilon_{hk} = \delta_{jh}\varepsilon_{ik}\) 且 \(1 = \sum_{i=1}^{n} \varepsilon_{ii}\)。证明 B 中存在元素族 \((e_{ij})\)（\(1 \leqslant i \leqslant n, 1 \leqslant j \leqslant n\)），满足 \(e_{ij}e_{hk} = \delta_{jh}e_{ik}\) 且 \(1 = \sum_{i=1}^{n} e_{ii}\)，并且对于每个指标有序对，\(\varepsilon_{ij}\) 是 \(e_{ij}\) 的典范像。（使用 (b)、代数第 VIII 章第 6 节的习题 11 以及代数第 VIII 章第 1 节的习题 9。）推出若 \(\overline{B}\) 同构于矩阵环 \(\mathbf{M}_{r}(\overline{\mathbf{D}})\)，其中 \(\overline{D}\) 是不必交换的域，则 B 同构于矩阵环 \(\mathbf{M}_{r}(\mathbf{D})\)，其中 D 是有限生成 A-模的 A-代数，且其 Jacobson 根 \(\mathfrak{d}\) 满足 \(D/\mathfrak{d}\) 同构于 \(\overline{D}\)（参见代数，第 VIII 章，第 1 节，习题 9 和 3）。证明若进一步 B 是 A 上的 Azumaya 代数（第二章，第 5 节，习题 14），则 D 也是。~
\end{enumerate}

\begin{enumerate}
\def\labelenumi{\arabic{enumi}.}
\setcounter{enumi}{5}
\tightlist
\item
  给出一个非交换 Artin 环 A 的例子，它带有由双边理想序列 \((a_{n})\) 构成的滤过，满足 \(a_{0} = A\)、\(a_{2} = 0\)，且第 6 节命题 8 对其不成立（参见代数，第 VIII 章，第 2 节，习题 6）。~
\end{enumerate}

¶ 7. (a) 设 A 为交换环，m 为 A 的理想，且 A 关于 m-进拓扑是 Hausdorff 且完备的，\(m/m^{2}\) 是有限生成 A-模。证明 \(A' = A\{X_{1}, \ldots, X_{r}\}\) 上的拓扑是 \(m'\)-进拓扑，其中 \(m' = mA'\)，且 \(m'/m'^{2}\) 是有限生成的 \((A'/m')\)-模（使用第 2 节，第 11 节，命题 14）。特别地，若 A 是 Noether 环，则 \(A'\) 也是 Noether 环。

\begin{enumerate}
\def\labelenumi{(\alph{enumi})}
\setcounter{enumi}{1}
\tightlist
\item
  设 A 为交换 Noether 环，m 为 A 的理想，且 A 关于 m-进拓扑是 Hausdorff 且完备的。令 \(u: A \rightarrow B\) 为从 A 到 Hausdorff 交换拓扑环 B 的连续同态，使 B 成为 A-代数。证明下列条件等价：
\end{enumerate}

(α) B 是 Noether 环，其拓扑是 mB-进拓扑，B 是完备的，且 B/mB 是 A/m 上的有限生成代数。

(β) B 作为拓扑 A-代数同构于 \(\lim B_{n}\)，其中 \((\mathrm{B}_{n})_{n\geqslant1}\) 是离散 A-代数的逆系统，满足当 \(\phi_{nm}:B_{m}\to B_{n}\) 时映射 \(m\geqslant n\) 是满射，\(\phi_{nm}\) 的核为 \(m^{n+1}B_{m}\)，且 \(B_{1}\) 是 A/m 上的有限生成代数。

(\((\gamma)\)) B 作为拓扑 A-代数同构于形如 \(\mathbf{A}\{\mathbf{X}_1,\dots ,\mathbf{X}_r\}\) 的代数对一个闭理想的商。

（为证明 \((\beta)\) 蕴含 \((\gamma)\)，使用第 2 节，第 11 节，命题 14 以及第 2 节，第 8 节，定理 1。）

\begin{enumerate}
\def\labelenumi{\arabic{enumi}.}
\setcounter{enumi}{7}
\tightlist
\item
  设 A 为带线性拓扑的交换环。加法群 \(A[[X_{1},\ldots,X_{p}]]\) 被识别为乘积群 \(A^{N^{p}}\)，并赋予乘积拓扑 T。
\end{enumerate}

\begin{enumerate}
\def\labelenumi{(\alph{enumi})}
\item
  证明 \(\mathcal{T}\) 与 \(\mathbf{A}[[\mathbf{X}_1,\dots ,\mathbf{X}_p]]\) 的环结构相容，并且是使映射 \(f\mapsto \partial f / \partial \mathbf{X}_i\) 连续的线性拓扑。
\item
  对于 \(P = \sum_{e} \alpha_{e,P} X^{e}\) 的每个元素 \(A[[X_{1}, \ldots, X_{p}]]\)（第 1 节，第 6 节的记号）以及每个完备 Hausdorff 拓扑 A-代数 B（第 2 节，习题 28），若元素 \(\mathbf{x} = (x_{1}, \ldots, x_{p}) \in B^{p}\) 满足族 \((\alpha_{e,P}\mathbf{x}^{e})\)（这里对 \(x^{e} = x_{1}^{e_{1}} \ldots x_{p}^{e_{p}}\) 写作 \(e = (e_{1}, \ldots, e_{p})\)）关于 \(N^{p}\) 的有限子集补集滤子在 B 中收敛到 0，则称 x 可代入 P。此族于是可求和，其和记作 \(P(\mathbf{x})\)。证明使 x 可代入 P 的形式幂级数 \(P \in A[[X_{1}, \ldots, X_{p}]]\) 的集合是 \(S_{x}\) 的子环 \(A[[X_{1}, \ldots, X_{p}]]\)，且映射 \(P \mapsto P(\mathbf{x})\) 是从 \(S_{x}\) 到 B 的同态。
\item
  证明若 \(\mathbf{x}\) 可代入 \(\mathbf{P}\)，且 \(\mathbf{y} = (y_1, \ldots, y_p)\) 是 \(\mathbf{B}^p\) 中的元素，并且对 \(y_i\)，\(1 \leqslant i \leqslant p\) 都是拓扑幂零的，则 \(\mathbf{x} + \mathbf{y}\) 可代入 \(\mathbf{P}\)。
\end{enumerate}

特别地，若 B 中存在一个由拓扑幂零元素组成的 0 的邻域，则对于所有 \(P \in A[[X_{1}, \ldots, X_{p}]]\)，可代入 P 的 \(x \in B^{p}\) 的集合 D 是开的，且从 D 到 B 的映射 \(\mathbf{x} \mapsto \mathrm{P}(\mathbf{x})\) 是连续的。

\begin{enumerate}
\def\labelenumi{(\alph{enumi})}
\setcounter{enumi}{3}
\tightlist
\item
  从现在起假设 A 是 Hausdorff 且完备的，并给 \(A[[X_{1},\ldots,X_{p}]]\) 赋予拓扑 T。为了使 \(P_{1},\ldots,P_{q}\) 中的 q 个形式幂级数 \(A[[X_{1},\ldots,X_{p}]]\) 可代入形式幂级数
\end{enumerate}

\[
\mathrm{Q} \in \mathrm{A} [ [ \mathrm{X} _ {1}, \dots , \mathrm{X} _ {q} ] ],
\]

其充要条件是系统 \((\mathbf{P}_{1}(0),\ldots,\mathbf{P}_{q}(0))\) 可代入 Q。

\begin{enumerate}
\def\labelenumi{(\alph{enumi})}
\setcounter{enumi}{4}
\item
  假设 \(\mathbf{x}\) 可代入每个 \(\mathrm{P}_k\)（\(1 \leqslant k \leqslant q\)），且系统 \((\mathrm{P}_1, \ldots, \mathrm{P}_q)\) 可代入 \(\mathbb{Q}\)。证明 \(\mathbf{x}\) 可代入 \(\mathbb{Q}(\mathrm{P}_1, \ldots, \mathrm{P}_q)\)，系统 \((\mathrm{P}_1(\mathbf{x}), \ldots, \mathrm{P}_q(\mathbf{x}))\) 可代入 \(\mathbb{Q}\)，并且当进一步满足下列假设之一时，\(\mathbb{Q}(\mathrm{P}_1(\mathbf{x}), \ldots, \mathrm{P}_q(\mathbf{x})) = ((\mathbb{Q}(\mathrm{P}_1, \ldots, \mathrm{P}_q))(\mathbf{x})\)：( \(\alpha\) ) \(\mathbb{B}\) 中存在理想 \(n\)，使环对 \((\mathbb{B}, n)\) 满足 Hensel 条件，且 \(x_i\)（\(1 \leqslant i \leqslant p\)）属于 \(n\)；( \(\beta\) ) 形式幂级数 \(\mathbb{Q}\) 是受限的。
\item
  取带离散拓扑的 \(\mathbf{B} = \mathbf{A} = \mathbf{F}_2\)，\(\mathbf{P} = \mathbf{X} + \mathbf{X}^2\)、\(\mathbf{Q} = \sum_{n=0}^{\infty} \mathbf{X}^{2^n}\)；则 \(\mathbf{P}\) 可代入 \(\mathbf{Q}\)，\(x = 1\) 可代入 \(\mathbf{P}\) 和 \(\mathbf{Q}(\mathbf{P})\)，\(\mathbf{P}(x)\) 可代入 \(\mathbf{Q}\)，但 \(\mathbf{Q}(\mathbf{P}(x)) = 0\)，而 \((\mathbf{Q}(\mathbf{P}))(x) = 1\)。
\end{enumerate}
